\documentclass[11pt,oneside,openany]{book}
\usepackage[T1]{fontenc}
\usepackage{lmodern}
\usepackage[letterpaper,margin=1.10in,includeheadfoot]{geometry}
\usepackage{amsmath,amssymb,amsthm,mathtools}
\usepackage{microtype,xcolor,titlesec,fancyhdr,etoolbox,needspace}
\usepackage{hyperref}

\ifdefined\pdftrailerid
  \pdftrailerid{MathGenome Topics in Representation Theory}
\fi
\definecolor{ink}{HTML}{172923}
\definecolor{deepgreen}{HTML}{2B5242}
\definecolor{rulegreen}{HTML}{8A9F92}
\hypersetup{colorlinks=true,linkcolor=deepgreen,urlcolor=deepgreen,citecolor=deepgreen,
  pdftitle={Topics in Representation Theory},pdfauthor={MathGenome},
  pdfsubject={ALG-07-D01 through ALG-07-D05},bookmarksnumbered=true}
\setcounter{tocdepth}{1}
\setlength{\headheight}{15pt}
\setlength{\emergencystretch}{2em}
\newcommand{\runningchapter}{}
\raggedbottom
\titleformat{\chapter}[display]
  {\normalfont\color{ink}}{\color{deepgreen}\large\scshape Chapter \thechapter}{8pt}
  {\huge\bfseries}
\titlespacing*{\chapter}{0pt}{7pt}{22pt}
\titleformat{\section}{\normalfont\Large\bfseries\color{ink}}{\thesection}{0.7em}{}
\pagestyle{fancy}
\fancyhf{}
\fancyhead[L]{\small\sffamily\color{deepgreen}MathGenome}
\fancyhead[R]{\small\nouppercase{\runningchapter}}
\fancyfoot[C]{\small\thepage}
\renewcommand{\headrulewidth}{0.35pt}
\renewcommand{\chaptermark}[1]{\renewcommand{\runningchapter}{\chaptername\ \thechapter}}
\fancypagestyle{plain}{\fancyhf{}\fancyfoot[C]{\small\thepage}
  \renewcommand{\headrulewidth}{0pt}}

\newtheorem{theorem}{Theorem}[chapter]
\newtheorem{lemma}[theorem]{Lemma}
\newtheorem{proposition}[theorem]{Proposition}
\newtheorem{corollary}[theorem]{Corollary}
\theoremstyle{definition}
\newtheorem{definition}[theorem]{Definition}
\newtheorem{example}[theorem]{Example}
\theoremstyle{remark}
\newtheorem{remark}[theorem]{Remark}
\newcommand*\keeptheoremheading[1]{\BeforeBeginEnvironment{#1}{\Needspace{4\baselineskip}}}
\forcsvlist{\keeptheoremheading}{theorem,lemma,proposition,corollary,definition,example,remark}
\allowdisplaybreaks[1]

\begin{document}
\frontmatter
\hypersetup{pageanchor=false}
\begin{titlepage}
\thispagestyle{empty}
\thispdfpagelabel{Cover}
\begin{flushleft}
\vspace*{0.72in}
{\large\scshape\color{deepgreen}MathGenome\par}
\vspace{0.36in}
{\color{rulegreen}\rule{0.94\textwidth}{0.65pt}\par}
\vspace{0.46in}
{\color{ink}\fontsize{35}{41}\selectfont\bfseries
  Topics in\\[2pt]Representation Theory\par}
\vspace{0.34in}
{\color{deepgreen}\Large Modular representations, quivers, invariants,\\[4pt]
  highest weights, and geometric Satake\par}
\vfill
{\normalsize ALG-07-D01 through ALG-07-D05\par}
\vspace*{0.50in}
\end{flushleft}
\end{titlepage}
\hypersetup{pageanchor=true}
\begingroup
\small
\makeatletter
\patchcmd{\l@chapter}{\vskip 1.0em \@plus \p@}
  {\vskip 0.45em \@plus 0.25\p@}{}{}
\makeatother
\tableofcontents
\endgroup

\chapter*{Scope and conventions}
\phantomsection
\addcontentsline{toc}{chapter}{Scope and conventions}
\renewcommand{\runningchapter}{Scope and conventions}

These notes continue the treatment of representations in
\emph{Algebra Core}, \emph{Rings, Modules, and Associative Algebras},
and \emph{The Symmetric Group and Schur--Weyl Duality}.
The chapters treat five different uses of a representation category.
For finite groups in characteristic \(p\), representations are
compared with characteristic zero through Brauer characters and
block idempotents. For quivers, reflection at sinks and sources
reduces the finite-type classification to an integral quadratic form.

The third chapter studies invariant rings and the effect of an
invariant projection on finite generation. It also gives proofs that
apply in modular examples where the projection can fail to exist.
The fourth chapter works with highest-weight categories, moving from
finite-dimensional semisimple Lie algebras to their untwisted affine
counterparts. The last chapter proves geometric Satake for complex
split tori with rational coefficients and studies positive-dimensional
convolution geometry for \(GL_2\).

Algebra modules are left modules unless stated otherwise. In the
modular chapter the splitting hypotheses will be specified separately
for the ordinary and modular coefficient fields. The proofs cite
earlier MathGenome results where their hypotheses match the present
argument. The bibliography provides references for each topic, and
the accompanying map records the location and proof dependencies of
each numbered statement.

\mainmatter
\chapter{Modular Characters and Blocks}
\label{ch:modular-representations}

Reduction from characteristic zero can turn an irreducible
representation into a module with several composition factors.
The Brauer character records those factors by evaluating on
elements of order prime to the characteristic. The block
decomposition then groups the simple modules and isolates the
local \(p\)-subgroups that contribute to each block.

We use the Jordan--H\"older theorem, the nilpotence of the
radical of a finite-dimensional algebra, and the Wedderburn--Artin
decomposition from \emph{Algebra Core}, Theorems~2.8, 3.3 and~3.4
\cite[PDF pp.~11, 14--15]{mg-algebra-core}. Maschke's theorem and
Schur's lemma are Theorems~8.2--8.3 there
\cite[PDF pp.~38--39]{mg-algebra-core}.
For Sylow existence, containment and conjugacy, we use D07
\cite[PDF pp.~49--50]{mg-d07}.

\section{Splitting systems and reduction}

\begin{definition}[A splitting \(p\)-modular system]
\label{res:MR01}
Let \(G\) be finite and \(p\) prime. Put \(m=|G|_{p'}\), the
largest divisor of \(|G|\) prime to \(p\). A splitting system for
the discussion consists of a complete discrete valuation ring
\(O\) with uniformizer \(\pi\), fraction field \(K\) of
characteristic zero and residue field \(k=O/\pi O\) of
characteristic \(p\). We assume that \(K\) contains the \(m\)-th
roots of unity and that the group algebras are split in the
following sense: \(KG\) and \(kG/J(kG)\) are finite products
of matrix algebras over their respective ground fields.
An element of \(G\) whose order is prime to \(p\) is called
\emph{\(p\)-regular}. We write \(G_{p'}\) for the set of such
elements.
\end{definition}

\begin{lemma}[Roots of unity and their residual lifts]
\label{res:MR02}
Reduction gives a group isomorphism from the \(m\)-th roots of
unity in \(O\) to those in \(k\). For a root \(\lambda\) in
\(k\), write \(\widehat{\lambda}\) for its unique lift.
If \(g\in G_{p'}\), an action of \(g\) on a finite-dimensional
\(k\)-representation is diagonalizable over \(k\) and its
eigenvalues are in this group of roots.
\end{lemma}

\begin{proof}
Every \(m\)-th root in \(K\) has valuation zero. If such a root
\(u\) reduces to \(1\), then
\[
 (u-1)(1+u+\dots+u^{m-1})=0 ,
\]
and the second factor is a unit, as its reduction is \(m\).
Hence \(u=1\). Applying this to the ratio of two roots shows
that the reduction of the \(m\) roots in \(K\) is injective.
Their images already account for every root of \(X^m-1\)
in \(k\). The operator given by \(g\) satisfies the polynomial
\(X^{|g|}-1\), which splits with distinct roots in \(k\).
\end{proof}

\begin{definition}[Brauer characters and ordinary characters]
\label{res:MR03}
For a finite-dimensional \(kG\)-module \(M\) and \(g\in G_{p'}\)
with eigenvalues \(\lambda_1,\ldots,\lambda_d\) on \(M\) counted
with multiplicities, its Brauer character is
\[
 \varphi_M(g)=\sum_{i=1}^{d}\widehat{\lambda_i}\in O.
\]
For a finite-dimensional \(KG\)-module \(V\), its ordinary
character \(\chi_V\) has value
\(\operatorname{tr}(g\mid V)\) on every \(g\in G\).
\end{definition}

\begin{lemma}[Additivity and reduction of a stable lattice]
\label{res:MR04}
Brauer characters are class functions on \(G_{p'}\) and are
additive on short exact sequences of \(kG\)-modules.
Every finite-dimensional \(KG\)-module \(V\) has a \(G\)-stable
\(O\)-lattice \(L\), where a lattice is a finite free submodule
spanning \(V\) over \(K\). For every such \(L\),
\[
 \varphi_{L/\pi L}=\chi_V|_{G_{p'}}.
\]
\end{lemma}

\begin{proof}
Conjugation of the matrix of \(g\) preserves its eigenvalues.
For a submodule the matrix of \(g\), in a basis starting with
a basis of that submodule, is block triangular. Its
characteristic polynomial is the product for the submodule
and the quotient. This proves additivity of the lifted sums.
For a lattice, take the \(O\)-span \(L_0\) of a \(K\)-basis and
put \(L=\sum_{g\in G}gL_0\). It is stable and spans \(V\).
After multiplication by a large enough power of \(\pi\), it is
a submodule of \(L_0\cong O^{\dim V}\). Such a submodule is free
by induction on the number of coordinates: its first-coordinate
image is a principal ideal, a lift of a generator of a nonzero
image splits the projection, and the kernel lies in the
remaining coordinates.

For \(g\in G_{p'}\) the polynomial \(X^{|g|}-1\) also diagonalizes
its action over \(K\). The characteristic polynomial in an
\(O\)-basis of \(L\) reduces to the characteristic polynomial
of the action on \(L/\pi L\). By Lemma~\ref{res:MR02}, the
eigenvalues in \(K\) all lie in \(O\), and distinct ones reduce
to distinct elements. Therefore each multiplicity is preserved
and the resulting lifted trace is the ordinary trace.
\end{proof}

\begin{theorem}[Brauer characters of simple modules form a basis]
\label{res:MR05}
Let \(S_1,\ldots,S_r\) be representatives of the simple
\(kG\)-modules up to isomorphism. The functions \(\varphi_{S_i}\)
are a \(K\)-basis of the space of \(K\)-valued class functions on
\(G_{p'}\). In particular \(r\) is the number of \(p\)-regular
conjugacy classes of \(G\).
\end{theorem}

\begin{proof}
Let \(t_i:kG\to k\) be the \(k\)-linear trace on \(S_i\).
These linear forms are independent. Indeed the quotient
\(kG/J(kG)\) is a product of full matrix algebras over \(k\),
with one factor acting on each \(S_i\). A lift of a matrix unit
in position \((1,1)\) in one factor has trace \(1\) on that
simple and trace \(0\) on every other simple.

For \(g\in G\), factor \(g=g_p g_{p'}\) inside the cyclic
subgroup generated by \(g\), with \(g_p\) of \(p\)-power
order and \(g_{p'}\) of order prime to \(p\). This follows by
choosing the two complementary idempotent exponents modulo
the order of \(g\) by the Chinese remainder theorem
\cite[PDF p.~66]{mg-d07}.
On the eigenspaces of \(g_{p'}\) in \(S_i\), the operator
\(g_p\) is unipotent, since \(X^{p^a}-1=(X-1)^{p^a}\)
in characteristic \(p\). Its trace on each eigenspace is the
dimension of that eigenspace, interpreted in \(k\). It follows
that \(t_i(g)=t_i(g_{p'})\). Hence a \(k\)-linear relation
among the restrictions \(t_i|_{G_{p'}}\) would give a relation
on the whole basis \(G\) of \(kG\) and then contradict independence.

If the \(\varphi_{S_i}\) satisfy a nonzero \(K\)-linear
relation, multiply the coefficients by a power of \(\pi\)
so that they all belong to \(O\) and at least one is a unit.
Reduction gives the corresponding nonzero relation among the
\(t_i|_{G_{p'}}\), a contradiction.

For spanning, the ordinary irreducible characters form a basis
of the class functions on \(G\) over \(K\). Indeed the split
matrix factors of \(KG\) give their independence by the trace
of matrix units, and the dimension of the center is the
number of factors. The center also has a basis of conjugacy
class sums, so that number is exactly the dimension of the
space of class functions. Every class function on \(G_{p'}\)
extends to \(G\) by setting its value on the other classes to
zero. Lemma~\ref{res:MR04} and a composition series of each
reduced lattice express the restricted ordinary characters
as linear combinations of the \(\varphi_{S_i}\).
\end{proof}

\begin{corollary}[The decomposition matrix]
\label{res:MR06}
For irreducible \(KG\)-modules \(V_\chi\), indexed by their
ordinary characters, let \(d_{\chi i}\) be the multiplicity
of \(S_i\) in a stable lattice for \(V_\chi\) reduced modulo
\(\pi\). The integers \(d_{\chi i}\) are independent of
the lattice and satisfy
\[
  \chi|_{G_{p'}}=\sum_{i=1}^{r} d_{\chi i}\varphi_{S_i}.
\]
The resulting map on the groups generated by module classes
with short-exact-sequence relations is additive.
\end{corollary}

\begin{proof}
Lemma~\ref{res:MR04} and additivity along a composition series
give the displayed formula. Its coefficients are uniquely
determined by Theorem~\ref{res:MR05}. The same argument works
for any finite-dimensional \(KG\)-module. Additivity follows
from the ordinary trace on short exact sequences and the
uniqueness of the coefficients on restriction.
\end{proof}

\section{Integral idempotents and projectives}

\begin{lemma}[Lifting an idempotent]
\label{res:MR07}
Suppose \(I\) is a nilpotent ideal of a unital ring \(R\), or
that \(R\) is complete and separated in its \(I\)-adic
topology. Every idempotent in \(R/I\) lifts to \(R\). If \(R\)
is commutative, the lifting is unique.
\end{lemma}

\begin{proof}
For a lift \(x\) put \(f=x^2-x\in I\). Then
\((2x-1)^2=1+4f\) is a unit by the finite or convergent
geometric series, so \(2x-1\) is a unit. All these elements
and inverses commute with \(x\). Replace \(x\) by
\(x'=x-(2x-1)^{-1}f\). Expanding the quadratic gives
\[
 x'^2-x'=(2x-1)^{-2}f^2 .
\]
Each step squares the power of \(I\) to which the error
belongs. The iteration either gives an idempotent in finitely
many steps or converges to one in the complete ring. In the
commutative case, two idempotents \(e,e'\) with the same
reduction give idempotents \(e(1-e')\) and \(e'(1-e)\) in
\(I\). An idempotent in \(I\) lies in every positive power
of \(I\), hence is zero in either case. Thus \(e=e'\).
\end{proof}

\begin{proposition}[Blocks and compatibility with reduction]
\label{res:MR08}
A block idempotent of \(kG\) is a primitive central
idempotent \(\bar b\). The block algebra is \(kG\bar b\).
The block idempotents are orthogonal, sum to \(1\), and
each simple module belongs to exactly one block, the one
whose idempotent acts as \(1\). Reduction gives a bijection
from the primitive central idempotents of \(OG\) to the block
idempotents of \(kG\). If \(V_\chi\) is an ordinary simple
module on which the lift \(b\) acts as \(1\), then
\(d_{\chi i}\) can be nonzero only for \(S_i\) in the block
\(\bar b\).
\end{proposition}

\begin{proof}
The finite-dimensional commutative algebra \(Z(kG)\) is a
product of local Artinian algebras by the commutative part of
the Artinian radical theorem cited above. The identities of
these factors are precisely the primitive central idempotents.
On a simple module every central idempotent has image zero or
the whole module, proving the assertion about its block.

Over \(O\) and over \(k\), the center of a group algebra has
the conjugacy class sums as a basis. Reduction on the center
is thus onto and has kernel \(\pi Z(OG)\); the latter center
is complete and separated for the \(\pi\)-adic topology.
Lemma~\ref{res:MR07} gives unique central lifts. Orthogonality
follows by lifting products, and a primitive central
idempotent on either side corresponds to one on the other
side, since a decomposition into orthogonal central
idempotents lifts or reduces. The only idempotent that
reduces to zero is zero. The lifted family again sums to \(1\).

If \(b\) acts as \(1\) on \(V_\chi\), then it acts as \(1\)
on every stable lattice inside it: \(b\) acts on the lattice
since \(b\in OG\). Therefore \(\bar b\) acts as \(1\) on the
whole reduction and on all its composition factors. This
gives the required zero entries of the decomposition matrix.
\end{proof}

\begin{theorem}[Brauer reciprocity and the Cartan matrix]
\label{res:MR09}
For \(A=kG\) choose an idempotent \(\bar e_i\in A\) lifting a
rank-one matrix idempotent in the \(S_i\)-factor of \(A/J(A)\),
with zero entries in the other factors. Put \(P_i=A\bar e_i\).
Then \(P_i\) is a projective cover of \(S_i\). For an
idempotent lift \(e_i\in OG\) of \(\bar e_i\), the
multiplicity of \(V_\chi\) in \(KGe_i\) is \(d_{\chi i}\).
Writing \(c_{ij}=[P_j:S_i]\) and \(D=(d_{\chi i})\), the
Cartan matrix is \(C=(c_{ij})=D^{\mathsf T}D\).
\end{theorem}

\begin{proof}
The lifts exist by Lemma~\ref{res:MR07}, first for the
nilpotent ideal \(J(A)\) and then for the ideal \(\pi OG\).
The left ideal \(P_i\) is a direct summand of the left
regular module and is projective. Its quotient by
\(J(A)P_i\) is the column module for the indicated simple
matrix factor. If a submodule \(U\subseteq P_i\) surjects
onto that column module, then \(P_i=U+J(A)P_i\), hence
\(P_i/U=J(A)(P_i/U)\). Iterating and using nilpotence
gives \(U=P_i\). This is the defining minimality of a
projective cover.

For any left \(A\)-module \(M\), evaluation at \(\bar e_i\)
identifies \(\operatorname{Hom}_A(P_i,M)\) with
\(\bar e_i M\) as \(k\)-vector spaces. The inverse sends
\(v\in\bar e_i M\) to \(a\bar e_i\mapsto av\).
The functor is exact by projectivity, and its value on
\(S_j\) has dimension \(1\) if \(j=i\) and \(0\) otherwise.
Hence \(\dim_k\bar e_i M=[M:S_i]\) for a finite-dimensional
module \(M\). Similarly, for a split semisimple \(KG\)-module
\(V\) one has
\(\operatorname{Hom}_{KG}(KGe_i,V)\cong e_iV\).
In particular, \(\dim_K e_i V_\chi\) is the multiplicity of
\(V_\chi\) in \(KGe_i\) by the split endomorphism ring of
each simple \(KG\)-module.

Let \(L\) be a stable lattice in \(V_\chi\).
Since \(e_iL\) is a direct summand of \(L\) as an
\(O\)-module, its reduction is \(\bar e_i(L/\pi L)\).
The rank of \(e_iL\) is both \(\dim_K e_iV_\chi\)
and \(\dim_k\bar e_i(L/\pi L)=d_{\chi i}\). This proves
the multiplicity assertion. Finally \(OGe_j\) is a stable
lattice in \(KGe_j\) whose reduction is \(P_j\).
By the multiplicity assertion the ordinary character of
\(KGe_j\) is \(\sum_\chi d_{\chi j}\chi\). Using
Corollary~\ref{res:MR06} on each summand and comparing
coefficients of \(\varphi_{S_i}\) gives
\(c_{ij}=\sum_\chi d_{\chi i}d_{\chi j}\).
\end{proof}

\section{The Brauer homomorphism and defect groups}

Write \(A^P\) for the elements of \(A=kG\) fixed under
conjugation by a subgroup \(P\). For \(Q\leq P\) define the
relative trace \(\operatorname{Tr}_Q^P:A^Q\to A^P\) by
\[
 \operatorname{Tr}_Q^P(a)=\sum_{u\in P/Q}uau^{-1},
\]
where the notation in the sum means a choice of representatives
of left cosets. Each summand depends only on the coset.

\begin{proposition}[The Brauer homomorphism]
\label{res:MR10}
If \(P\leq G\) is a \(p\)-subgroup, the coefficient projection
\[
 \operatorname{Br}_P:A^P\longrightarrow kC_G(P),\qquad
 \sum_{x\in G}a_xx\longmapsto\sum_{x\in C_G(P)}a_xx
\]
is a surjective algebra homomorphism. Its kernel is the sum
of \(\operatorname{Tr}_Q^P(A^Q)\) over all proper subgroups
\(Q<P\). It takes central elements of \(kG\) to central
elements of \(kC_G(P)\) and commutes with transport by
conjugation in \(G\).
\end{proposition}

\begin{proof}
For \(z\in C_G(P)\) the pairs \(x,y\in G\) with \(xy=z\) form
a \(P\)-set under simultaneous conjugation. If \(a,b\in A^P\),
the product \(a_xb_y\) of their coefficients is constant on
each orbit. Every orbit of size greater than one contributes
zero in characteristic \(p\). A pair in a singleton orbit
has \(x,y\in C_G(P)\), so the coefficient of \(z\) in \(ab\)
is exactly the coefficient in
\(\operatorname{Br}_P(a)\operatorname{Br}_P(b)\).
The projection is the identity on the subalgebra \(kC_G(P)\),
so it is unital and surjective.

An invariant element is a linear combination of \(P\)-orbit
sums. The singleton orbit sums are exactly the surviving basis
elements under projection. The orbit sum of any other \(x\in G\)
is \(\operatorname{Tr}_{C_P(x)}^P(x)\), from the proper
stabilizer of \(x\). Conversely, each proper-subgroup trace
has zero coefficient at every \(z\in C_G(P)\) because that
coefficient is a multiple of \([P:Q]\). This proves the kernel
formula. Conjugating the coefficients commutes with projection
if one conjugates \(P\) at the same time. In particular every
element of \(C_G(P)\) fixes the image of a central element.
\end{proof}

\begin{lemma}[Relative trace ideals in the center]
\label{res:MR11}
For a \(p\)-subgroup \(P\leq G\), put
\(I_P=\operatorname{Tr}_P^G(A^P)\subseteq Z(A)\).
It is an ideal of \(Z(A)\). For the class sum \(\kappa_x\)
of the conjugacy class of \(x\), one has \(\kappa_x\in I_P\)
if and only if a Sylow \(p\)-subgroup of \(C_G(x)\) is
conjugate in \(G\) to a subgroup of \(P\). Such class sums
form a basis of \(I_P\).
\end{lemma}

\begin{proof}
Multiplication by a central element commutes with the
relative trace, and the product in the domain is still
\(P\)-invariant, so \(I_P\) is an ideal. On the \(P\)-orbit
sum of \(x\) the trace is
\[
 \operatorname{Tr}_{P}^{G}
    \bigl(\operatorname{Tr}_{C_P(x)}^P(x)\bigr)
  = \operatorname{Tr}_{C_P(x)}^G(x)
  = [C_G(x):C_P(x)]\kappa_x .
\]
The first equality follows by multiplying coset
representatives. The last equality counts the cosets of
\(C_P(x)\) over each coset of \(C_G(x)\).
The integer scalar is nonzero in \(k\) exactly when
\(C_P(x)\) is a Sylow \(p\)-subgroup of \(C_G(x)\).
An orbit sum with this property occurs in the conjugacy
class of \(x\) exactly when a conjugate of such a Sylow
subgroup is contained in \(P\). Since the class sums are
linearly independent, the trace images of the orbit-sum
basis of \(A^P\) establish both assertions.
\end{proof}

\begin{theorem}[Defect groups and their conjugacy]
\label{res:MR12}
For each block idempotent \(b\in kG\) there exists a
\(p\)-subgroup \(D\leq G\) such that for every \(p\)-subgroup
\(P\leq G\),
\[
 \operatorname{Br}_P(b)\ne0
 \quad\Longleftrightarrow\quad
 P\text{ is contained in a conjugate of }D .
\]
In particular the inclusion-maximal \(P\) for which
\(\operatorname{Br}_P(b)\ne0\) form one conjugacy class,
called the defect groups of \(b\). A representative can
be obtained as a Sylow \(p\)-subgroup of a centralizer
\(C_G(x)\) of an element with a nonzero coefficient in \(b\).
It has largest possible order among the Sylow subgroups of
the centralizers of such elements.
\end{theorem}

\begin{proof}
Write \(b=\sum_C a_C\kappa_C\) as a sum over conjugacy classes,
restricting to \(a_C\ne0\). For each such \(C\) choose
\(x_C\in C\) and a Sylow \(p\)-subgroup \(D_C\) of \(C_G(x_C)\).
By Lemma~\ref{res:MR11}, \(\kappa_C\in I_{D_C}\).
The ring \(b Z(A)\) is local, by the primitive central block
decomposition in Proposition~\ref{res:MR08}. In this ring
\[
 b=b^2=\sum_C b a_C\kappa_C .
\]
At least one term on the right is a unit of \(bZ(A)\),
since otherwise their sum belongs to its maximal ideal.
For its class \(C\), multiplication by the inverse in \(bZ(A)\)
shows \(b\in I_{D_C}\), because \(I_{D_C}\) is a central
ideal. Put \(D=D_C\).

Every class in the support of \(b\) now has a centralizer
whose Sylow \(p\)-subgroups are conjugate into \(D\), again
by the class-sum description of \(I_D\).
If \(\operatorname{Br}_P(b)\ne0\), some element in the support
of \(b\) centralizes \(P\). Then \(P\) is contained in a Sylow
\(p\)-subgroup of its centralizer, hence in a conjugate of
\(D\). Conversely \(x_C\) centralizes \(D\) and occurs in
the support of \(b\). For any \(P\) contained in a conjugate
of \(D\), a conjugate of \(x_C\) shows
\(\operatorname{Br}_P(b)\ne0\). This also proves the
maximality and order statements.
\end{proof}

\begin{corollary}[Principal blocks and defect zero]
\label{res:MR13}
The principal block, defined to be the block of the trivial
representation, has as defect groups the Sylow \(p\)-subgroups
of \(G\). A block has trivial defect groups if and only if
its block algebra is semisimple.
\end{corollary}

\begin{proof}
On the trivial representation a group algebra element acts by
its augmentation, the sum of its coefficients. Thus the
augmentation of the principal idempotent is \(1\).
For an invariant element, projection \(\operatorname{Br}_P\)
does not change its augmentation, since the dropped
\(P\)-orbit sums contribute multiples of \(p\).
For \(P\) Sylow in \(G\) this proves that the principal block
has \(\operatorname{Br}_P(b)\ne0\).
Theorem~\ref{res:MR12} gives the claimed defect groups.

For the second assertion, Theorem~\ref{res:MR12} and
Lemma~\ref{res:MR11} show that trivial defect is equivalent
to \(b\in I_1\). If \(b=\operatorname{Tr}_1^G(a)\) and \(M\)
is a finite-dimensional module in the block, consider the
surjection \(A\otimes_k M\to M\), \(x\otimes m\mapsto xm\),
where \(G\) acts on the first factor only. It has the
\(G\)-linear section
\[
 m\longmapsto\sum_{g\in G} g\otimes a(g^{-1}m).
\]
Equivariance is a reindexing, and the composite is
\(\operatorname{Tr}_1^G(a)\) acting on \(M\), namely the
identity. Thus every such \(M\) is projective over \(A\).
Splitting the simple quotients in a composition series of the
left regular block module shows that the block is semisimple.

Conversely, a split semisimple block \(bA\) consists of one
matrix factor, since \(b\) is primitive central. Its simple
column \(M\) is a summand of \(bA\), hence projective over \(A\).
Choose an \(A\)-linear section \(s\) of the same surjection
\(A\otimes_k M\to M\), and let \(f:M\to M\) be the
coefficient at the basis element \(1\in G\) of \(s\).
Equivariance makes the coefficient at \(g\) equal to
\(f(g^{-1}m)\). Hence
\(\sum_{g\in G}g f(g^{-1}m)=m\). The block \(bA\) acts on
its column as the full endomorphism algebra, so \(f\) is
the action of some \(a\in bA\). The central element
\(\operatorname{Tr}_1^G(a)\) belongs to \(bA\) and acts as
\(1\) on its column. The faithful column action then gives
\(\operatorname{Tr}_1^G(a)=b\), so \(b\in I_1\).
\end{proof}

\section{An example at two different primes}

\begin{example}[The symmetric group of degree three]
\label{res:MR14}
Over splitting systems as in Definition~\ref{res:MR01},
the decomposition matrices and block defects for \(S_3\)
are as follows. The ordinary rows in both matrices are
the trivial, sign, and two-dimensional standard representations.
At \(p=2\) the columns are the trivial and the simple
two-dimensional representation. At \(p=3\) the columns
are the trivial and sign representations.
\[
 D_2=\begin{pmatrix}1&0\\1&0\\0&1\end{pmatrix},
 \qquad
 D_3=\begin{pmatrix}1&0\\0&1\\1&1\end{pmatrix}.
\]
The respective Cartan matrices are
\(\operatorname{diag}(2,1)\) and
\(\left(\begin{smallmatrix}2&1\\1&2\end{smallmatrix}\right)\).
At \(p=2\) there are two blocks, the principal block with
Sylow defect and the other block with trivial defect.
At \(p=3\) the principal block is the only block, with its
Sylow subgroup of order three as defect group.
\end{example}

\begin{proof}
Write \(c=(1\,2\,3)\) and \(t=(1\,2)\). The standard
representation \(W\) in characteristic zero is the
coordinate-sum-zero subspace of the permutation representation.
Over a splitting field for \(c\) it has the two distinct
eigenlines with nontrivial cube-root eigenvalues. They are
interchanged by \(t\), since \(tct^{-1}=c^{-1}\).
So \(W\) is absolutely irreducible. The three indicated
ordinary simples account for the full algebra dimension
\(1^2+1^2+2^2=6\). The traces on \(W\) at \(1,t,c\) are
respectively \(2,0,-1\), by subtracting the trace of the
complementary trivial representation from that of the permutation
representation. The group-stable lattice in \(W\) given by
the coordinate-sum-zero \(O\)-module reduces to the
coordinate-sum-zero subspace of \(k^3\): the defining linear map
to \(O\) is a split surjection of \(O\)-modules.

In characteristic two this reduced \(W\) is simple, by the
same two-eigenline argument. Those lines are defined over \(k\)
by Lemma~\ref{res:MR02}. The \(2\)-regular classes are \(1\)
and the class of \(c\), so Theorem~\ref{res:MR05} shows that
the trivial module and \(W\) are the two simples. The restrictions
of the ordinary trivial and sign characters coincide.
The character identity of Corollary~\ref{res:MR06} gives
the displayed \(D_2\). The class sum \(b=c+c^{-1}\) is central,
and \(b^2=b\) in characteristic two. It acts as zero on the
trivial module and as \(1\) on \(W\), since the sum of the
nontrivial cube roots is \(-1=1\) in \(k\).
Every nonzero central idempotent acts nontrivially on at least
one simple module: otherwise it belongs to the nilpotent
radical. Thus both \(b\) and \(1-b\) are primitive central,
and \(1-b\) is the principal idempotent.
The two elements in the support of \(b\) have centralizer
of order three, so Theorem~\ref{res:MR12} gives trivial
defect for \(b\). Corollary~\ref{res:MR13} gives the
defect of \(1-b\).

In characteristic three the \(3\)-regular classes are \(1\)
and the class of \(t\). Trivial and sign are distinct simples
and exhaust the simple modules by Theorem~\ref{res:MR05}.
Their Brauer values at \(t\) are respectively \(1,-1\).
Comparing the ordinary values at the two \(3\)-regular
classes gives \(D_3\). In the reduction of \(W\), the
constant-coordinate line is a trivial submodule, since a
constant vector now has coordinate sum zero. The quotient
has dimension one and is the sign module by the two
composition factors found in \(D_3\).
This extension does not split: both one-dimensional modules
restrict trivially to \(\langle c\rangle\), whereas \(c\)
acts nontrivially on the reduced \(W\), for example on the
vector \((1,-1,0)\). A short exact sequence between modules
in distinct blocks splits by applying the orthogonal central
idempotents, so trivial and sign belong to the same block.
There can be no other block, by the nontrivial action of
a nonzero central idempotent on some simple. The Sylow
defect follows as above.
The Cartan matrices follow from Theorem~\ref{res:MR09}.
\end{proof}

\paragraph{Sources.}
For the theory of stable lattices and Brauer characters see
Webb \cite[Sections~9.4--9.5 and 10.1--10.3]{webb} and
Serre \cite[Chapters~15--16 and 18]{serre-finrep}.
For blocks, the relative-trace approach to defect, and the
Brauer morphism, see Webb \cite[Sections~12.1--12.5]{webb}.
\chapter{Quivers of Finite Representation Type}

We retain the path convention of \emph{Rings, Modules, and Associative
Algebras}, RM03-C06.02--05~\cite[PDF pp.~14--15]{mg-rm03}: an
arrow \(i\to j\) acts from the \(i\)-space to the \(j\)-space in a
left \(kQ\)-module. The quiver \(Q\) is finite and has no oriented
cycles. Representations in this chapter are finite dimensional.
The proof of Gabriel's theorem uses linear reflection functors and
an integral quadratic form, so it keeps track of parallel edges
and of arbitrary acyclic orientations.

\section{An Euler form from the arrows}

\begin{definition}[The Euler form and the Tits form]
\label{res:QU01}
For dimension vectors \(d,e\in\mathbb Z^{Q_0}\), put
\[
 \langle d,e\rangle=\sum_i d_i e_i-\sum_{a:i\to j}d_i e_j,
 \qquad q(d)=\langle d,d\rangle.
\]
Write \(a_{ij}\) for the number of edges between \(i\) and \(j\),
independent of their orientation. The symmetric matrix of the
polarization of \(q\) has diagonal entries \(2\) and off-diagonal
entries \(-a_{ij}\). Write \(e_i\) for the \(i\)-th unit vector
when discussing dimension vectors and \(S_i\) for the representation
with a copy of \(k\) at \(i\) and zero at the other vertices.
\end{definition}

\begin{proposition}[The two-term projective resolution]
\label{res:QU02}
For \(P_i=kQ e_i\) and a representation \(V\) there is an exact
sequence of modules
\[
 0\longrightarrow \bigoplus_{a:i\to j}P_j\otimes_k V_i
   \xrightarrow{\partial}\bigoplus_i P_i\otimes_k V_i
   \longrightarrow V\longrightarrow0,
 \qquad
 \partial(p\otimes v)=pa\otimes v-p\otimes V_a v .
\]
The final map sends a path generator \(p\otimes v\) to \(V_p v\).
Consequently, writing \(d=\dim V\) and \(e=\dim W\), one has
\[
 \dim_k\operatorname{Hom}_{kQ}(V,W)
   -\dim_k\operatorname{Ext}^1_{kQ}(V,W)=\langle d,e\rangle,
 \qquad \operatorname{Ext}^r_{kQ}(V,W)=0\quad(r\ge2).
\]
\end{proposition}
\begin{proof}
By the path basis of RM03-C06.03, \(P_i\) has the paths starting
at \(i\) as basis. The idempotents split \(kQ\) as the direct sum
of the \(P_i\), so they are projective. The displayed maps are
module maps and compose to zero. Every nontrivial path has a unique
first arrow \(a\), followed by a path \(p\), and thus occurs in
exactly one first term of \(\partial\). If a nonzero element on
the left has a largest path length among its nonzero coefficients,
the first terms of that length plus one in its image cannot cancel.
This proves injectivity. In the cokernel the same relation reduces
the length of paths, down to \(e_i\otimes v\). These remaining
vectors map isomorphically to the direct sum of the vertex spaces.
This proves exactness. The correspondence and vertexwise
exactness are those of RM03-C06.05~\cite[PDF pp.~14--15]{mg-rm03}.

Apply \(\operatorname{Hom}_{kQ}(-,W)\) to this resolution.
Evaluation at \(e_i\) identifies
\(\operatorname{Hom}_{kQ}(P_i\otimes V_j,W)\) with
\(\operatorname{Hom}_k(V_j,W_i)\). The resulting two-term complex
computes \(\operatorname{Ext}\), by the projective-resolution
formula in Result R11 of \emph{Advanced Homological Algebra}
~\cite[PDF p.~32]{mg-ha}. Its alternating dimension gives the
claimed formula.
\end{proof}

\section{Reflection at a sink or a source}

\begin{proposition}[The reflection equivalence]
\label{res:QU03}
At a sink \(i\), let \(m:\bigoplus_{a:j\to i}V_j\to V_i\)
be the sum of the incoming arrow maps. Reverse the incident arrows
and replace \(V_i\) by \(\ker m\). The new arrows are the
projections out of this kernel. Dually, at a source replace the
source space by the cokernel of the map to the sum of the
outgoing target spaces. The new arrows are the inclusions followed
by that quotient. These constructions, leaving the other arrows
unchanged, are inverse equivalences on the full subcategories
of representations with no summand \(S_i\).

For a representation in those subcategories the new dimension
vector is \(s_i(d)\), where
\[
 (s_i d)_j=d_j\quad(j\ne i),\qquad
 (s_i d)_i=-d_i+\sum_j a_{ij}d_j.
\]
The map \(s_i\) is an involution preserving \(q\).
\end{proposition}
\begin{proof}
At a sink, a complementary subspace to the image of \(m\) in
\(V_i\) is a direct summand supported at \(i\). Thus having no
\(S_i\)-summand means \(m\) is onto. At a source, a complementary
subspace to the kernel of the outgoing map at \(i\) shows that
the analogous condition is injectivity. Reflection at a sink
produces an injective outgoing map. The next reflection at the
same vertex takes its cokernel and recovers \(V_i\) and its
incoming maps from the exact sequence of \(m\). Starting at a
source uses the dual exact sequence. The identifications commute
with all morphisms, since vertex maps preserve the respective
kernels and induce maps on the cokernels. A direct summand of
an object with no \(S_i\) again has no \(S_i\), so these
equivalences preserve indecomposability. Rank-nullity gives the
dimension rule. As a polynomial in \(d_i\), the form \(q\) is
\[
 d_i^2-d_i\sum_j a_{ij}d_j+q(d\text{ with its \(i\)-entry set to }0).
\]
Substitution of the new entry leaves this expression unchanged
and applying the substitution twice recovers \(d_i\).
\end{proof}

\begin{lemma}[A full cycle of sink reflections]
\label{res:QU04}
Number the vertices \(1,\ldots,n\) so that all arrows point from
higher to lower indices. Reflect in the order \(1,\ldots,n\).
Each step takes place at a sink, and the quiver at the end has
its original orientation. Put \(C=s_n\cdots s_1\) and let \(E\)
be lower triangular with diagonal \(1\) and \((i,j)\)-entry
\(-a_{ij}\) for \(i>j\). Then \(C=-E^{-1}E^{\mathsf T}\).
If \(q\) is positive definite, \(C\) has finite order and no
nonzero fixed vector. For every nonzero \(x\) with nonnegative
coordinates, some power of \(C\) has a negative coordinate on \(x\).
\end{lemma}
\begin{proof}
Acyclicity supplies the numbering by successively taking a sink
of the remaining quiver. At stage \(i\), the edges to smaller
indices have reversed and now point toward \(i\). The other
edges still point toward \(i\). Every edge is reversed twice in
the full cycle. If \(y=Cx\), at the \(i\)-th update the earlier
entries are already \(y_j\) and the later ones are still \(x_j\).
Thus
\[
  y_i-\sum_{j<i}a_{ij}y_j=-x_i+\sum_{j>i}a_{ij}x_j .
\]
This is \(Ey=-E^{\mathsf T}x\). Since \(E+E^{\mathsf T}\) is
the positive definite matrix of the polarization, \(C-I\)
is invertible over \(\mathbb R\). There are only finitely many
integer vectors having any fixed \(q\)-value, since \(q\) bounds
the Euclidean norm from below by a positive multiple of its
square. Hence there are only finitely many integral invertible
linear maps preserving \(q\), by considering the images of the
unit vectors. In particular \(C^h=I\) for some \(h>0\).
Multiplying the geometric sum by \(C-I\) shows that
\(\sum_{j=0}^{h-1}C^j x=0\). This would be impossible if all its
summands were entrywise nonnegative and \(x\ne 0\).
\end{proof}

\section{The form and finite representation type}

\begin{proposition}[The positive definite graphs]
\label{res:QU05}
The form \(q\) is positive definite exactly when every connected
component of the underlying graph is \(A_n,D_n,E_6,E_7\), or
\(E_8\). Here \(D_n\) has \(n\ge4\), and an isolated vertex is \(A_1\).
\end{proposition}
\begin{proof}
An edge of multiplicity at least two gives \(q\le0\) on the vector
with value \(1\) at its endpoints and zero elsewhere. A cycle
gives \(q\le0\) by assigning \(1\) on its vertices. A vertex of
degree at least four gives \(q\le0\) by putting \(2\) there and
\(1\) on four of its neighbors. In a tree with two trivalent
vertices, put \(2\) along the path joining them and \(1\) at their
four adjacent vertices off the path. If there are \(m\) path
vertices, the sum of squares is \(4m+4\), which equals the sum of
edge products \(4(m-1)+8\) on this support. Thus only paths and
trees with a single trivalent vertex can be positive definite.

On a path, twice the form is the sum of squares of successive
differences after adjoining zero at both ends. On a tree with one
trivalent vertex, let the arm lengths be \(p-1,r-1,s-1\) in
edges and put \(x\) at the center. Attach a zero beyond each
leaf. The same difference-square expression on the three arms,
less \(x^2\), equals \(2q\), by expanding at the center where
three squares meet. For an arm with \(p\) differences their
sum is \(x\), up to sign. The minimum of their squared sum is
\(x^2/p\), attained by the equally spaced values. Thus for
\(x\ne0\) the least possible \(2q\) is
\(x^2(1/p+1/r+1/s-1)\). For \(x=0\) every nonzero vector gives
a positive value. With \(2\le p\le r\le s\), the strict inequality
\(1/p+1/r+1/s>1\) forces \(p=2\). If \(r=2\) one obtains
\((2,2,s)\), the \(D_{s+2}\) family. If \(r=3\) one needs
\(s=3,4,5\), which gives \(E_6,E_7,E_8\). If \(r\ge4\) the
sum is at most \(1\). These computations prove the classification.
\end{proof}

\begin{lemma}[A dimension obstruction to finite type]
\label{res:QU06}
Over an infinite field \(k\), if \(q\) is not positive definite
then \(Q\) has infinitely many isomorphism classes of
indecomposable finite-dimensional representations.
\end{lemma}
\begin{proof}
There is a nonzero \(d\in\mathbb Z_{\ge0}^{Q_0}\) with \(q(d)\le0\).
If \(q\) has negative values, choose a rational vector with
negative value by continuity. If it is semidefinite and singular,
the rational polarization matrix has a nontrivial rational
kernel. In either case replace coordinates by their absolute
values, which cannot increase \(q\), and clear denominators.

For this \(d\), the arrow matrices form \(k^N\) with
\(N=\sum_{a:i\to j}d_i d_j\). Isomorphism is the change-of-basis
action of \(\prod_i GL_{d_i}(k)\), with \(D=\sum_i d_i^2\)
matrix entries before imposing invertibility. Multiplying all
basis matrices by the same scalar has no effect. Choose one
vertex of positive dimension and cover the action by the finitely
many choices of a nonzero entry of its basis matrix. Normalize
that entry to \(1\) using the common scalar. In each such chart,
the \(N\) arrow entries on the orbit of a fixed representation
are rational functions of at most \(D-1<N\) variables, since
\(q(d)=D-N\le0\). They are algebraically dependent over \(k\),
by the transcendence-degree bound of D07
~\cite[Theorem 21.35, PDF p.~176]{mg-d07}.
A nonzero polynomial relation between them vanishes everywhere on
the corresponding orbit chart. The product of these finitely
many nonzero relations vanishes on the entire orbit.

If there were only finitely many orbits, their polynomial
relations would give a nonzero polynomial vanishing on all of
\(k^N\). A polynomial over an infinite field vanishing at every
point is zero, by induction on the number of variables and the
root bound in one variable. This contradicts the relation.
Every representation splits as a sum of indecomposables by
induction on total dimension. If there were only finitely many
indecomposables, then at a fixed total dimension only finitely
many finite lists of them could occur. Thus there would be
only finitely many orbits for \(d\), a contradiction.
\end{proof}

\begin{theorem}[Gabriel's theorem and the dimension-vector bijection]
\label{res:QU07}
Over an algebraically closed field \(k\), a finite acyclic quiver
has finite representation type exactly when its underlying graph
is a disjoint union of the Dynkin graphs \(A,D,E\).
In that case the dimension-vector map is a bijection from
isomorphism classes of indecomposables to the positive roots of
the graph. In coordinates these are exactly the nonnegative
integral vectors \(d\) with \(q(d)=1\).
Every such indecomposable has endomorphism ring \(k\)
and has no nonzero self-extension class.
\end{theorem}
\begin{proof}
Suppose \(q\) is positive definite. Apply the cyclic sequence of
sink reflections from Lemma~\ref{res:QU04} to an indecomposable.
At a sink the process continues unless the representation is
\(S_i\), by Proposition~\ref{res:QU03}. If it never stopped,
every power of \(C\) would send its original dimension vector
to a nonnegative vector. Lemma~\ref{res:QU04} rules this out.
The indecomposable therefore reduces to a vertex simple, and
preservation of \(q\) gives \(q(d)=1\). The equivalences also
identify its endomorphism ring with that of the simple, namely
\(k\). The self-extension assertion follows from
Proposition~\ref{res:QU02}.

Conversely take a nonnegative integral \(d\) with \(q(d)=1\).
As long as its entries stay nonnegative, apply the same sink
reflections to the vector. There is a first step producing a
negative coordinate by Lemma~\ref{res:QU04}. Just before that
step, the entry at its vertex \(i\) is strictly larger than
\(\sum_j a_{ij}d_j\). The expression for \(q\) in
Proposition~\ref{res:QU03} shows that the contribution from
that coordinate is at least \(d_i\), and the remaining form is
positive if any remaining entry is nonzero. Since \(q(d)=1\),
the vector at that step is \(e_i\). Start with \(S_i\) at
that stage and apply the inverse source reflections in reverse
order. None of these inverse steps is obstructed: a source vector
\(e_j\) at a step would reflect to \(-e_j\), contradicting the
nonnegativity of the recorded previous vector. The functors
therefore produce the required indecomposable. If two
indecomposables have dimension \(d\), the common sequence ends
at the same simple and gives an isomorphism back at the
original orientation.

Write \(W=\langle s_i\rangle\) and define the graph's root set
as the union of the orbits \(W e_i\). These vectors span and
have square length \(2\) for the polarized form. Reflection in
\(w e_i\) is \(w s_i w^{-1}\), so it permutes the set. Integral
pairings and equal lengths give the crystallographic and reduced
conditions. The set is finite by positive definiteness. If an
integral vector \(x\) has both signs, write \(x=x^+-x^-\) with
disjoint nonnegative supports. The off-diagonal entries of the
form give \(q(x)\ge q(x^+)+q(x^-)\ge2\). Thus these roots are
positive or negative in the simple-vector basis. Every positive
root has \(q=1\), and the preceding vector reduction expresses
every nonnegative vector of \(q\)-value \(1\) as such an orbit point.
There are finitely many such vectors by positive definiteness.
Proposition~\ref{res:QU05} now gives the asserted types, and
Lemma~\ref{res:QU06} proves the converse of finite type because
an algebraically closed field is infinite.
\end{proof}

\begin{example}[The almost split sequence for a single arrow]
\label{res:QU08}
For \(1\to2\), the indecomposables are \(S_1,S_2\), and \(P_1\),
where \(P_1\) has identity arrow between two copies of \(k\).
The sequence
\[
 0\longrightarrow S_2\longrightarrow P_1
   \longrightarrow S_1\longrightarrow0
\]
is almost split.
\end{example}
\begin{proof}
The three vectors of \(q\)-value \(1\) are \((1,0),(0,1),(1,1)\),
as predicted by Theorem~\ref{res:QU07}. Over any field, the kernel,
image and cokernel of a single linear map give a direct sum
decomposition into copies of these three objects. The displayed sequence uses the
inclusion at vertex \(2\) and the projection at vertex \(1\).
It is exact at each vertex and nonsplit, since there is no
nonzero map \(S_1\to P_1\) compatible with the arrow.
For a map from a direct sum of indecomposables to \(S_1\) that is
not a split epimorphism, the components from \(S_1\) vanish
and those from \(S_2\) are zero. The remaining components
from \(P_1\) are scalar multiples of the displayed projection,
so the map factors through it. Dually, a map out of \(S_2\)
that is not a split monomorphism has zero components to \(S_2\)
and to \(S_1\). Each component to \(P_1\) is a scalar
multiple of the inclusion, giving the other factorization.
These are the two almost split properties. The general existence
theorem for finite-dimensional algebras, with its nonprojective
or noninjective endpoint hypotheses, is Result M14 in
\emph{Topics in Rings, Modules, and Associative Algebras}
~\cite[PDF p.~19]{mg-rm03d}.
\end{proof}

\paragraph{Sources.}
For the Euler form and the reflection proof, see
Etingof et al.~\cite[Sections 5.4--5.8]{etingof-rep}.
For the path algebra and the general theory of almost split
sequences see Assem--Simson--Skowro\'nski~\cite[Chapters II and IV]{assem}.
\chapter{Invariant Rings and Reynolds Operators}

Let \(k\) be a field. An action on a \(k\)-algebra acts by
\(k\)-algebra automorphisms. For a smooth algebraic group over an
algebraically closed field, we use rational actions and identify
invariance with invariance under all of its \(k\)-points. A
rational action is locally finite: for a vector \(v\), a
coaction identity \(g(v)=\sum_j f_j(g)v_j\) with finitely many
terms places its orbit in a finite-dimensional subspace, and
the span of that orbit is stable. Finite abstract groups also
act locally finitely. We use the Hilbert basis theorem in D07,
Result K07.22.hilbert\_basis~\cite[Section 23.11, PDF p.~191]{mg-d07}.

\section{Invariant projection and exactness}

\begin{definition}[Invariants and Reynolds operators]
\label{res:IN01}
The invariants of a \(G\)-module \(M\) are \(M^G=\{v:gv=v
\text{ for every }g\in G\}\). For a \(G\)-algebra \(A\), this is
a subalgebra. A Reynolds operator on \(A\) is a \(G\)-equivariant
\(A^G\)-linear map \(R:A\to A^G\) whose restriction to \(A^G\)
is the identity. An action category is called linearly reductive
when every finite-dimensional representation in it is a direct
sum of simple representations.
\end{definition}

\begin{proposition}[Exact invariants and the canonical projection]
\label{res:IN02}
If the action category is linearly reductive, every locally finite
representation \(M\) has a natural projection \(M\to M^G\). It
kills every simple submodule with nontrivial action. On an algebra
this projection is a Reynolds operator. Invariants are exact on
locally finite representations in this case. Conversely, exactness
on finite-dimensional representations implies complete reducibility.
\end{proposition}
\begin{proof}
In a finite-dimensional semisimple representation take the sum of
the trivial simple submodules and the sum of the other simple
submodules. They form a direct sum, since a map from a nontrivial
simple to a trivial one is zero, and the first sum equals the
invariant subspace. Their projections are preserved by equivariant
maps: a map of simple modules has image zero or an isomorphic
simple. Any invariant subspace of finite dimension is contained
in this system of compatible projections. On a locally finite
module, use the projection of a finite-dimensional stable
subspace containing the given vector. Two choices fit in their
sum, so the projection is well defined. Multiplication by
\(a\in A^G\) is equivariant, so naturality gives \(R(ax)=aR(x)\).

For an equivariant surjection, project any lift of an invariant
vector to obtain an invariant lift. Kernels of the invariants
map are invariants in the kernel by the definition of fixed
vectors. Conversely for a finite-dimensional equivariant
surjection \(M\to N\), the map
\(\operatorname{Hom}_k(N,M)\to\operatorname{Hom}_k(N,N)\)
is onto by elementary linear algebra. It is equivariant for
the conjugate action on the Hom spaces, and in the rational
case the Hom spaces are again rational. Exactness on invariants
lifts the identity of \(N\) to an equivariant section. Every
short exact sequence of finite-dimensional representations
therefore splits, yielding complete reducibility by induction
on dimension.
\end{proof}

\begin{proposition}[Averaging for finite groups and split tori]
\label{res:IN03}
For a finite group \(G\) with \(|G|\) invertible in \(k\), the
operator \(R(v)=|G|^{-1}\sum_{g\in G}gv\) gives exact invariants
and the Reynolds operator of Proposition~\ref{res:IN02}.
For a split torus \(\mathbb G_m^r\) acting rationally, projection
to weight zero gives exact invariants in every characteristic.
\end{proposition}
\begin{proof}
In the finite-group case, reindexing the sum gives equivariance
and invariance, and the sum fixes an invariant vector. It is
preserved by equivariant maps, so applying it to any lift of an
invariant vector proves exactness. Proposition~\ref{res:IN02}
then identifies it with the canonical projection.

For a torus, write the coaction of \(m\) as
\(\sum_{\lambda\in\mathbb Z^r}m_\lambda\otimes z^\lambda\),
a finite sum of Laurent monomials. Coassociativity and the
identities
\(\Delta z^\lambda=z^\lambda\otimes z^\lambda\) show that
each \(m_\lambda\) has weight \(\lambda\), by independence of
the Laurent monomials. The counit gives \(m=\sum_\lambda m_\lambda\).
Applying the coaction again makes the sum direct. Equivariant
maps preserve each weight, which proves exactness and the
claim about the projection. The same argument works over a
finite \(k\) if invariants are taken for the algebraic group
coaction, since weight zero is then the defining condition.
\end{proof}

\section{Finite generation}

\begin{theorem}[Hilbert finite generation using Reynolds]
\label{res:IN04}
Let \(A=\bigoplus_{d\ge0} A_d\) be a finitely generated
graded commutative \(k\)-algebra with \(A_0=k\).
Suppose \(G\) preserves the grading and there is a
degree-preserving Reynolds operator on \(A\). Then \(A^G\)
is generated by finitely many homogeneous invariants.
\end{theorem}
\begin{proof}
By the Hilbert basis theorem the ideal generated in \(A\)
by all positive-degree homogeneous invariants is generated by
finitely many of them, say \(f_1,\ldots,f_m\). Indeed any
finite ideal generating set is made of sums using only
finitely many of the original generators. For a homogeneous
invariant \(f\) of positive degree, write \(f=\sum_i a_i f_i\).
Taking components of degree \(\deg f\) makes each nonzero
\(a_i\) homogeneous of lower degree. Apply \(R\) to obtain
\(f=\sum_i R(a_i)f_i\). Induction on degree gives
\(A^G=k[f_1,\ldots,f_m]\), since the fixed subring is the sum
of its homogeneous components.
\end{proof}

\begin{corollary}[Rational actions of characteristic-zero reductive groups]
\label{res:IN05}
Let \(k\) be algebraically closed of characteristic zero and
let \(G\) be a connected reductive algebraic group over \(k\).
For every finitely generated commutative rational \(G\)-algebra
\(A\), the invariant algebra \(A^G\) is finitely generated.
The construction also applies to a rational split torus
action in arbitrary characteristic over an algebraically
closed field.
\end{corollary}
\begin{proof}
Result E07 of \emph{Linear Algebraic Groups}
~\cite[PDF p.~40]{mg-alggroup} gives complete reducibility for
the first case. For the torus use Proposition~\ref{res:IN03}.
Choose a finite-dimensional stable subspace \(V\subset A\)
containing \(1\) and algebra generators. It exists by local
finiteness. Let \(A_d'\) be the span of products of \(d\)
elements of \(V\), with \(A_0'=k\). The spaces are nested
because \(1\in V\) and have union \(A\).
The algebra \(\bigoplus_{d\ge0}A_d' t^d\subseteq A[t]\)
is generated by \(Vt\) and is graded with degree-zero part \(k\).
The projection in Proposition~\ref{res:IN02} preserves each
stable homogeneous degree. Apply Theorem~\ref{res:IN04}.
For \(a\in A^G\) lying in \(A_d'\), the element \(at^d\) is
invariant, and evaluation at \(t=1\) recovers \(a\).
Evaluation therefore maps the finitely generated invariant
algebra of the graded construction onto \(A^G\).
\end{proof}

\begin{theorem}[Finite-group invariants in arbitrary characteristic]
\label{res:IN06}
If a finite group acts on a finitely generated commutative
\(k\)-algebra \(A\), then \(A^G\) is a finitely generated
\(k\)-algebra and \(A\) is finite as an \(A^G\)-module.
\end{theorem}
\begin{proof}
Choose algebra generators \(a_1,\ldots,a_n\). The coefficients
of the monic orbit polynomial
\(\prod_{g\in G}(T-g(a_i))\) are invariant. Let \(B\) be the
\(k\)-subalgebra generated by all these coefficients.
Each \(a_i\) satisfies its orbit polynomial, since the identity
factor vanishes at \(T=a_i\). Thus bounded powers of the
generators span \(A\) over \(B\): reduce high powers using those
monic equations. The ring \(B\) is noetherian, so the
\(B\)-submodule \(A^G\) is finite. Its module generators
together with the listed generators of \(B\) generate it as
a \(k\)-algebra. The same finite list that spans \(A\) over
\(B\) spans it over \(A^G\).
\end{proof}

\begin{theorem}[The Noether bound in characteristic zero or large characteristic]
\label{res:IN07}
Let \(G\) have order \(m\) and act linearly on the
degree-one space of \(S=k[x_1,\ldots,x_n]\). If \(k\) has
characteristic zero or characteristic greater than \(m\),
then \(S^G\) is generated by invariants of degrees at most \(m\).
\end{theorem}
\begin{proof}
Let \(I\subset S\) be generated by the homogeneous invariants of
degrees \(1,\ldots,m\). Take a formal linear combination
\(z=\sum_i t_i x_i\) using variables fixed by \(G\). The
coefficients in the \(t_i\) of the orbit polynomial coefficients
of \(z\) are invariant in \(S\), of degrees \(1,\ldots,m\).
Indeed each is an elementary symmetric sum of the \(g(z)\),
and the action permutes these terms. Substitution \(T=z\)
gives \(z^m\in I\,S[t_1,\ldots,t_n]\).
Each degree-\(m\) monomial \(x^\alpha\) therefore belongs
to \(I\), by comparing its coefficient of \(t^\alpha\)
and dividing by the multinomial coefficient
\(m!/\prod_i\alpha_i!\), which is invertible under the
stated hypotheses. Hence every homogeneous element of
degree at least \(m\) lies in \(I\). Choose bases of the
finitely many invariant spaces of degrees at most \(m\).
For a homogeneous invariant of larger degree express it
in \(I\) using these bases, and apply the averaging of
Proposition~\ref{res:IN03} to its coefficients.
Induction as in Theorem~\ref{res:IN04} proves the bound.
\end{proof}

\section{Computations with and without averaging}

\begin{proposition}[Molien's formula in characteristic zero]
\label{res:IN08}
For a finite group \(G\) acting on a finite-dimensional
space \(W\) over a field of characteristic zero,
\[
 \sum_{d\ge0} \dim_k(\operatorname{Sym}^d W)^G\,t^d
       =\frac{1}{|G|}\sum_{g\in G}\frac{1}{\det(1-tg|_W)}.
\]
The identity is one of formal power series over \(k\), using
the integer dimensions as coefficients.
\end{proposition}
\begin{proof}
The degree-\(d\) dimension equals the trace of the degree-\(d\)
averaging projection, since a projection is identity on its
image and zero on its kernel. For a fixed \(g\), extend scalars
for the calculation to an algebraic closure. Its finite order
in characteristic zero makes \(g\) diagonalizable, with
eigenvalues \(\lambda_1,\ldots,\lambda_n\), since its minimal
polynomial divides a separable polynomial \(X^{|G|}-1\).
The trace on degree \(d\) sums \(\prod_i\lambda_i^{a_i}\) over
\(\sum_i a_i=d\) by the monomial basis. Summing the geometric
series gives \(\prod_i(1-\lambda_i t)^{-1}\) and hence the
stated determinant. Average and use Proposition~\ref{res:IN03}.
\end{proof}

\begin{example}[A polynomial modular invariant ring without a Reynolds operator]
\label{res:IN09}
In characteristic \(p>0\), let a generator \(g\) of \(C_p\)
act on \(k[x,y]\) by \(g(x)=x\) and \(g(y)=y+x\).
Then the invariant ring is \(k[x,y^p-x^{p-1}y]\), and
there is no equivariant projection onto it.
\end{example}
\begin{proof}
The substitution has order \(p\) and fixes
\(N=y^p-x^{p-1}y\) by the binomial formula. The polynomial
ring is free over \(k[x,N]\) on \(1,y,\ldots,y^{p-1}\).
The identity \(y^p=N+x^{p-1}y\) spans it, and independence
follows from the distinct residues modulo \(p\) of the
largest \(y\)-degrees in \(b_j(x,N)y^j\) for nonzero
\(b_j\). If a fixed element has highest basis index
\(j>0\), applying \(g-1\) gives coefficient \(jxb_j\) at
\(y^{j-1}\), nonzero in the polynomial ring. Thus only
the basis index \(0\) is possible. An equivariant
projection \(P\) would satisfy
\(P(x)=P((g-1)y)=(g-1)P(y)=0\), although \(P\) fixes
the nonzero invariant \(x\).
\end{proof}

\begin{example}[Three vectors for the special linear group]
\label{res:IN10}
Let \(k\) be algebraically closed of any characteristic.
For the diagonal action of \(\operatorname{SL}_2(k)\)
on three column vectors \(v_1,v_2,v_3\in k^2\), the
invariant polynomial functions form the polynomial ring
on the three brackets \(d_{ij}=\det(v_i,v_j)\), \(i<j\).
\end{example}
\begin{proof}
The brackets are invariant by the determinant.
On the set \(d_{12}\ne0\), the matrix
\(\operatorname{diag}(1,d_{12})[v_1\ v_2]^{-1}\) has determinant
one. It takes the first two columns to \((1,0)^{\mathsf T}\)
and \((0,d_{12})^{\mathsf T}\). In these coordinates the
third is \((-d_{23}/d_{12},d_{13})^{\mathsf T}\).
These normalized triples realize all values of
\(d_{12}\in k^\times\) and \(d_{13},d_{23}\in k\).
This proves independence of the brackets, because a
polynomial vanishing for all these values is zero over
the infinite field \(k\). An invariant polynomial, evaluated
on the normalized triples, is a Laurent polynomial in
\(d_{12}\) with polynomial dependence on the other two brackets.
The equality on this open set is an equality of rational
functions: after clearing denominators, multiply by \(d_{12}\)
to get a polynomial vanishing at all triples, hence zero.

Interchanging columns \(2\) and \(3\) gives the same assertion
with denominator a power of \(d_{13}\). The intersection of the
two localizations of the polynomial ring on the brackets is
the polynomial ring itself. Indeed if
\(a/d_{12}^u=b/d_{13}^v\), then
\(a d_{13}^v=b d_{12}^u\); the two polynomial variables
are relatively prime and prime, so \(d_{12}^u\) divides \(a\).
This proves generation in every characteristic.
\end{proof}

\paragraph{Scope.}
The general positive-characteristic reductive-group finite
generation theorem uses geometric or power reductivity. Its
statement and the power-surjectivity hypotheses are developed
by van der Kallen~\cite[Theorem 8]{vdk}. For the Reynolds,
degree-bound, and bracket calculations, see
Derksen--Kemper~\cite[Sections 2.2, 3.1--3.2 and 4.5]{derksen}
and Kraft--Procesi~\cite[Sections 1.6 and 8.4]{kraft-procesi}.
\chapter[Category O and Affine Highest Weights]{Category \(\mathcal O\) and Affine Highest Weights}
\label{ch:category-o}

The finite-dimensional theory starts from a highest weight and asks
when its simple module is finite-dimensional. Category \(\mathcal O\)
keeps the Verma modules and their extensions, so that multiplicities
in these modules can be compared with filtrations of projective covers.
For an affine loop algebra the same triangular construction is available.
The central charge and the additional simple root then enter the
integrability condition.

\section{Weight spaces and finite length}

Work over \(\mathbb C\). In the finite-dimensional part,
\(\mathfrak g\) is semisimple with a triangular decomposition
\(\mathfrak n^-\oplus\mathfrak h\oplus\mathfrak n^+\), and
\(e_i,f_i,h_i\) are its simple-root triples. We use the triangular
PBW isomorphism and the universality and simple quotient of the
Verma module \(M(\lambda)\) from \cite[Results G01--G03, PDF
pp.~43--44]{mg-lie}. Write \(L(\lambda)\) for that simple quotient,
\(Q=\sum_i\mathbb Z\alpha_i\) for the root lattice,
\(Q_+=\sum_i\mathbb Z_{\ge0}\alpha_i\), and \(\nu\le\lambda\) if
\(\lambda-\nu\in Q_+\).

\begin{definition}[The category \(\mathcal O\)]
\label{res:LO01}
An object of \(\mathcal O\) is a finitely generated \(U(\mathfrak g)\)-module
which is a direct sum of its \(\mathfrak h\)-weight spaces and on which
\(\mathfrak n^+\) is locally finite. Local finiteness means that
\(U(\mathfrak n^+)v\) is finite-dimensional for every vector \(v\).
The morphisms are the \(U(\mathfrak g)\)-linear maps.
\end{definition}

\begin{theorem}[Bounds, simple objects and finite length]
\label{res:LO02}
Every object of \(\mathcal O\) has finite-dimensional weight spaces
and has weights contained in a finite union of sets \(\lambda-Q_+\).
The category is abelian, with kernels and cokernels as for modules,
and its objects have finite length. Its simple objects are exactly
the pairwise nonisomorphic \(L(\lambda)\).
\end{theorem}
\begin{proof}
In any weight module each \(\mathfrak h\)-stable subspace is the
sum of its intersections with the weight spaces. To see this,
separate the finitely many weights of a vector by an element of
\(\mathfrak h\) and project onto each component by polynomial
interpolation. Thus weight spaces are exact under subquotients.

Choose finitely many weight-vector generators of an object, by
replacing each generator by its weight components. The sum \(E\) of
their \(U(\mathfrak n^+)\)-spans is finite-dimensional and stable under
\(\mathfrak b=\mathfrak h\oplus\mathfrak n^+\). The module is a quotient of
\(U(\mathfrak g)\otimes_{U(\mathfrak b)}E\).
By PBW the induced module is \(U(\mathfrak n^-)\otimes E\) as a
weight space decomposition. For a fixed difference of weights, only
finitely many negative-root monomials have that weight
\cite[Result G02, PDF p.~44]{mg-lie}. This gives both asserted bounds.
A Verma module belongs to \(\mathcal O\) because it is cyclic,
has finite-dimensional weight spaces, and has only finitely many
weights reachable by upward shifts from any fixed weight.
A maximal weight exists in any nonempty subset of these supporting
cones: above a given weight in the subset the union of the cones
contains only finitely many weights. A vector at a maximal weight
is killed by the positive root operators. In a simple module it
generates, and the Verma universal property identifies the simple
module with \(L(\lambda)\). The highest weights distinguish the simples.

We prove finite length first for \(M(\lambda)\). Let \(Z\) be the
center of \(U(\mathfrak g)\), let \(W\) be the finite Weyl group, and
let \(\rho\) be half the sum of the positive roots. The center acts
on \(M(\lambda)\) by \(\chi_\lambda\), and
\(\chi_\nu=\chi_\lambda\) if and only if
\(\nu+\rho\in W(\lambda+\rho)\), by the Harish--Chandra theorem
\cite[Result H08, PDF pp.~54--55]{mg-lie}.
Every nonzero subquotient of \(M(\lambda)\) has a maximal weight
\(\nu\). A vector at that weight gives a nonzero map from
\(M(\nu)\), so the scalar central actions on source and target
agree. Set
\(S=W(\lambda+\rho)-\rho\). Every nonzero subquotient therefore
has positive sum of the dimensions of its weight spaces over \(S\).
This sum is exact on short exact sequences and is finite on
\(M(\lambda)\). It bounds the length of every strict chain of
submodules, proving finite length.

The finite weight module \(E\) above has a flag of
\(\mathfrak b\)-submodules with one-dimensional weight quotients.
Order its finitely many weights with larger weights first,
then take a flag in each weight space. Positive root operators
raise weights, so all terms are stable. Induction to
\(\mathfrak g\) is exact because the right \(U(\mathfrak b)\)-module
\(U(\mathfrak g)\) is free by PBW. It therefore turns the flag
into a filtration by Verma modules. The induced module and its
quotient have finite length.

A submodule of a finite-length \(U(\mathfrak g)\)-module is finitely
generated. Indeed it has finite length itself, and one can lift
one cyclic generator of each simple factor. The weight and local
finiteness properties pass to subquotients and finite direct sums.
Hence the module kernels and cokernels are in \(\mathcal O\),
with the usual identification of coimage and image.
\end{proof}

\begin{proposition}[Central character and root-lattice components]
\label{res:LO03}
Each object of \(\mathcal O\) decomposes functorially and exactly
into a finite direct sum indexed by generalized central characters
and cosets of \(Q\) in \(\mathfrak h^*\).
In the component for \((\chi,\xi)\), the simple labels are those
\(\mu\in\xi\) with \(\chi_\mu=\chi\). They form a finite set and
have a common upper bound \(\Lambda\in\xi\). Every weight of every
object in that component is at most \(\Lambda\).
\end{proposition}
\begin{proof}
Root operators preserve the coset of each weight, and the bounds
of Theorem~\ref{res:LO02} imply that only finitely many cosets occur.
In a fixed object there are finitely many distinct central
characters of simple factors. Choose \(z\in Z\) whose values on
them are pairwise distinct, by choosing a linear combination of
a finite set of separating central elements outside finitely many
proper hyperplanes. On an object of length \(l\), the operator
\(\prod_{\chi}(z-\chi(z))^l\) vanishes, since each factor for a
simple quotient kills that quotient. The corresponding coprime
polynomial factors give the generalized eigenspace decomposition,
by the polynomial Bezout identity. On each term the simple factors
all have the indicated character, so every \(z'\in Z\) acts there
by \(\chi(z')\) plus a nilpotent operator. This intrinsic
description proves functoriality. To verify exactness in a short
exact sequence, use one separating \(z\) and a common exponent
for all three objects. The resulting polynomial projectors commute
with the maps of the sequence.

The possible labels are in one finite shifted Weyl orbit by
\cite[Result H08, PDF pp.~54--55]{mg-lie}. In a fixed coset write
them in simple-root coordinates relative to any chosen weight of
the coset and take the coordinatewise maximum of their finitely
many integral differences. This gives a common \(\Lambda\).
The weights of a finite-length object are the union of the
weights of its simple factors, by exactness on weight spaces.
Each \(L(\mu)\) has weights at most \(\mu\), so the common bound
applies to the entire component.
\end{proof}

\section{Projective covers and reciprocity}

\begin{lemma}[A truncated weight projective]
\label{res:LO04}
For \(\mu\le\Lambda\), let \(\mathcal O^{\le\Lambda}\) be the
full subcategory of objects with all weights at most \(\Lambda\).
It has a projective object \(P_\Lambda(\mu)\) representing the
exact functor \(X\mapsto X_\mu\). This projective has a Verma
filtration whose successive quotients, grouped by their distinct
highest weights \(\nu\), are direct sums of \(M(\nu)\) with
\(\mu\le\nu\le\Lambda\).
\end{lemma}
\begin{proof}
Give \(U(\mathfrak n^+)\) a generator of weight \(\mu\): multiplication
is the \(\mathfrak n^+\)-action and elements of degree \(\beta\in Q_+\)
have weight \(\mu+\beta\). This defines an action of
\(\mathfrak b\) because the adjoint action of \(\mathfrak h\) on
the degree-\(\beta\) part is multiplication by \(\beta\).
Kill the direct sum of spaces of weights \(\nu\not\le\Lambda\).
This is a \(\mathfrak b\)-submodule because upward shifts keep a
weight outside the down-set of \(\Lambda\). Write \(E_\Lambda(\mu)\)
for the quotient. It is finite-dimensional: its allowed degrees
have simple-root coordinates between \(0\) and those of
\(\Lambda-\mu\), and each PBW-degree space for a fixed root-lattice
degree is finite-dimensional.

Define
\[
  P_\Lambda(\mu)=U(\mathfrak g)\otimes_{U(\mathfrak b)}E_\Lambda(\mu).
\]
The negative-root PBW basis gives the bound on its weights and
finite-dimensional weight spaces. A vector can be raised only
into finitely many weights at most \(\Lambda\); hence its
\(U(\mathfrak n^+)\)-span is finite-dimensional. The induced
module is cyclic and lies in \(\mathcal O^{\le\Lambda}\).
A vector \(x\in X_\mu\) in this subcategory gives a unique
\(\mathfrak b\)-map \(E_\Lambda(\mu)\to X\) taking the generator
to \(x\), since all the removed degrees act by zero there.
Induction extends it uniquely to a \(\mathfrak g\)-map. This
represents the weight functor and proves projectivity.

Order the distinct weights of \(E_\Lambda(\mu)\) with larger weights
first. The subspaces spanned by the initial weight spaces form
a \(\mathfrak b\)-stable filtration. Each quotient consists of
copies of \(\mathbb C_\nu\) for its weight \(\nu\). Exact induction
turns these quotients into copies of \(M(\nu)\), as asserted.
\end{proof}

\begin{theorem}[Projective covers have Verma filtrations]
\label{res:LO05}
For each \(\mu\) there is a projective cover
\(P(\mu)\twoheadrightarrow L(\mu)\) in \(\mathcal O\).
It has a finite filtration by Verma modules. In a filtration
obtained from the truncated projective, the displayed highest
weights can be arranged with larger weights first.
\end{theorem}
\begin{proof}
Take the component for \((\chi_\mu,\mu+Q)\) in
Proposition~\ref{res:LO03} and its bound \(\Lambda\).
Apply its exact projector to \(P_\Lambda(\mu)\). Denote the
result by \(P_0\). Any map from \(P_\Lambda(\mu)\) into an object
of that component factors through \(P_0\), so its projectivity
in the cutoff category gives projectivity there. The exact
component decomposition then makes \(P_0\) projective in
\(\mathcal O\). The map representing a nonzero vector of
\(L(\mu)_\mu\) gives a surjection from \(P_0\).

Decompose \(P_0\) as a finite direct sum of indecomposables by
induction on its length, and take a summand \(P\) whose map
\(p:P\to L(\mu)\) is nonzero. The map is surjective and its
kernel is superfluous, meaning that adding it to a proper
submodule never gives \(P\). In fact if a submodule \(J\) maps
onto \(L(\mu)\), projectivity lifts \(p\) to a map into \(J\).
Composing with the inclusion gives an endomorphism \(a\) of \(P\)
with \(pa=p\). In a finite-length module the sequences of kernels
and images of powers stabilize. For large \(n\) this gives
\(P=\ker a^n\oplus\operatorname{im}a^n\): the intersection is
zero by kernel stabilization, and for any \(x\) choose \(y\)
with \(a^nx=a^{2n}y\) by image stabilization.
Then \(x-a^ny\in\ker a^n\). Since \(P\) is indecomposable and \(pa^n=p\ne0\),
\(a\) is invertible. Thus \(J=P\). This proves the cover property,
and also uniqueness of the maximal submodule \(\ker p\).
Two projective covers of the same simple lift their maps through
one another. Each lift maps onto the target by superfluity,
and the finite lengths force the lifts to be isomorphisms.

It remains to transfer the Verma filtration to \(P\).
In the filtration of \(P_\Lambda(\mu)\) constructed in
Lemma~\ref{res:LO04}, the term indexed by the first \(i\) distinct
weights \(\nu_1,\ldots,\nu_i\) is the submodule generated by the
weight spaces at those weights in the full module. Indeed the
term was induced from these weight spaces in \(E_\Lambda(\mu)\),
and the remaining quotient has a Verma filtration with highest
weights all after the first \(i\) in the chosen order.
Its weights cannot equal any of \(\nu_1,\ldots,\nu_i\), since
they are at most their respective later highest weights and
the specified order places larger weights first. The generation
description persists under the central-character projector on
\(P_\Lambda(\mu)\), as it preserves weight spaces. Consequently
every endomorphism of \(P_0\) preserves the induced filtration.

Apply the projection onto \(P\). The \(i\)-th quotient of the
resulting filtration is the image of an idempotent on a direct
sum of copies of \(M(\nu_i)\), or zero. An endomorphism of one
copy is scalar because the generating top weight space is
one-dimensional. Thus idempotents on a sum of copies are
ordinary complex idempotent matrices and their images are
direct sums of the same Verma module. Refining these quotients
gives a Verma filtration of \(P\) in the stated order.
\end{proof}

\begin{lemma}[Dual Vermas detect Verma factors]
\label{res:LO06}
There is an exact contravariant involution \(D\) on
\(\mathcal O\) preserving each simple \(L(\lambda)\) up to
isomorphism. Write \(\nabla(\lambda)=D M(\lambda)\). Then
\[
 \dim_{\mathbb C}\operatorname{Hom}\bigl(M(\nu),\nabla(\lambda)\bigr)
   =\begin{cases}1&\nu=\lambda,\\0&\nu\ne\lambda,\end{cases}
 \qquad
 \operatorname{Ext}^1_{\mathcal O}\bigl(M(\nu),\nabla(\lambda)\bigr)=0 .
\]
\end{lemma}
\begin{proof}
The map fixing \(h_i\) and interchanging \(e_i\) and \(f_i\)
extends to an algebra anti-involution \(\tau\) on \(U(\mathfrak g)\).
For the Lie relations one reverses brackets on both sides:
the Cartan relations then transform into one another,
and \([e_i,f_j]=\delta_{ij}h_i\) becomes
\([e_j,f_i]=\delta_{ij}h_i\). A Serre relation with nested
positive brackets becomes the corresponding negative one
times an overall sign. These relations present \(\mathfrak g\)
by \cite[Results F01--F03, PDF pp.~36--38]{mg-lie}.
For a finite-weight-space module define \(DX\) as the direct
sum of the linear duals of its weight spaces, with action
\((u\ell)(x)=\ell(\tau(u)x)\). A homogeneous \(u\) shifts
the support of a functional by the weight of \(u\), so this
action preserves finite support. It is exact and involutive
on these weight modules by finite-dimensional duality and
exactness on each weight.

It sends a simple module to a simple one because applying
the exact double dual to a subquotient recovers that subquotient.
The maximal weight of \(DL(\lambda)\) is still \(\lambda\),
so \(DL(\lambda)\cong L(\lambda)\). For any object of
\(\mathcal O\), its dual has finite length by the dual of a
composition series, hence is finitely generated. Its weight
support is unchanged and finite-dimensional weight spaces
remain finite-dimensional. The upper bounds thus imply local
finiteness for \(\mathfrak n^+\), proving \(DX\in\mathcal O\).

A highest vector in \(\nabla(\lambda)\) is a functional
annihilating all the subspaces \(f_iM(\lambda)\).
Their sum contains every weight space of \(M(\lambda)\) below
the top: the \(f_i\) generate \(\mathfrak n^-\), and every
nonempty word in them starts with some \(f_i\).
A highest functional is therefore at the top weight \(\lambda\)
and its space is one-dimensional. The Hom formula follows
from the universal property of \(M(\nu)\).

For the Ext assertion take a short exact sequence
\(0\to\nabla(\lambda)\to X\to M(\nu)\to0\) in \(\mathcal O\).
If \(\nu\not<\lambda\), lift the highest vector at weight
\(\nu\). Raising this lift lands in \(\nabla(\lambda)\) at a
weight \(\nu+\alpha_i\), which is impossible because it would
imply \(\nu<\lambda\). Thus the lift is highest and furnishes
a splitting by universality. If \(\nu<\lambda\), dualize
to an extension of \(M(\lambda)\) by \(\nabla(\nu)\) and
apply the first case. Dualizing its splitting gives the
original splitting.
\end{proof}

\begin{theorem}[BGG reciprocity]
\label{res:LO07}
For a Verma filtration of \(P(\mu)\), the number of factors
isomorphic to \(M(\lambda)\) is independent of the filtration
and satisfies
\[
 [P(\mu):M(\lambda)]=[M(\lambda):L(\mu)] .
\]
The right-hand side is a composition multiplicity.
\end{theorem}
\begin{proof}
The unique maximal submodule of the cover from
Theorem~\ref{res:LO05} implies that every nonzero map to a
simple module factors through the top \(L(\mu)\).
Its endomorphism algebra is \(\mathbb C\), by the
one-dimensional generating highest vector. Therefore
\(\operatorname{Hom}(P(\mu),L(\eta))\) has dimension
\(\delta_{\mu,\eta}\). Projectivity and a composition
series then give
\(\dim\operatorname{Hom}(P(\mu),\nabla(\lambda))
=[\nabla(\lambda):L(\mu)]=[M(\lambda):L(\mu)]\),
the last equality by restricted duality.

Apply \(\operatorname{Hom}(-,\nabla(\lambda))\) to the successive
short exact sequences of a Verma filtration of \(P(\mu)\).
For a short exact sequence with Verma quotient, any map
from its submodule to \(\nabla(\lambda)\) extends: push
out along the map and split the resulting extension by
Lemma~\ref{res:LO06}. Thus the dimensions add across the
filtration. The Hom formula in the same lemma gives the
multiplicity \([P(\mu):M(\lambda)]\).
\end{proof}

\begin{example}[The principal rank-one component]
\label{res:LO08}
For \(\mathfrak g=\mathfrak{sl}_2\) write
\(\lambda=\lambda(h)\in\mathbb C\). The Verma module
\(M(\lambda)\) is simple unless \(\lambda=m\in\mathbb Z_{\ge0}\).
In that case its only nonzero proper submodule is the
Verma module \(M(-m-2)\), generated by \(f^{m+1}v_m\).
In the central-character component with labels \(0,-2\),
the projectives are \(P(0)=M(0)\) up to isomorphism and a
nonsplit extension \(0\to M(0)\to P(-2)\to M(-2)\to0\).
\end{example}
\begin{proof}
PBW gives the vectors \(f^j v_\lambda\) as a basis, with
weight \(\lambda-2j\) and \(j\ge0\). Commuting \(e\)
through one additional \(f\), using \([e,f]=h\) and
\([h,f]=-2f\), gives inductively
\[
  e f^j v_\lambda=j(\lambda-j+1) f^{j-1}v_\lambda .
\]
A nonzero proper submodule has a nonzero weight vector of
least exponent \(j>0\); applying \(e\) forces \(j=\lambda+1\).
Then every vector with exponent at least \(j\) is obtained
by powers of \(f\), so these span the whole submodule.
At \(m\ge0\) that span is a Verma module: its highest
weight is \(m-2(m+1)=-m-2\), and the universal map is
injective because the images of its PBW basis are distinct
nonzero basis vectors.
The Harish--Chandra theorem gives the labels \(0,-2\)
in this component. Reciprocity gives one Verma factor
for \(P(0)\), and both \(M(0)\) and \(M(-2)\) for \(P(-2)\).
The ordering in Theorem~\ref{res:LO05} places \(M(0)\) in
the submodule. A split extension would contradict the
indecomposability of the projective cover constructed there.
\end{proof}

\section{The untwisted affine loop algebra}

For the remainder of the chapter assume \(\mathfrak g\) is simple.
Fix its invariant nondegenerate symmetric form \(B\) normalized so
that long roots have squared length \(2\) for the induced form on
\(\mathfrak h^*\) \cite[Results B06 and D02--D05, PDF pp.~15,
23--25]{mg-lie}. Write \(x_m=x\otimes t^m\) for a loop vector.

\begin{proposition}[The loop cocycle and the degree derivation]
\label{res:LO09}
The vector space
\(\widehat{\mathfrak g}=
(\mathfrak g\otimes\mathbb C[t,t^{-1}])\oplus\mathbb C K\oplus\mathbb C d\)
is a Lie algebra for
\[
 [x_m,y_n]=[x,y]_{m+n}+m\,\delta_{m,-n}B(x,y)K,\qquad
 [d,x_m]=m x_m,\qquad [K,\widehat{\mathfrak g}]=0 .
\]
Its subalgebra spanned by the loops and \(K\) is a nonsplit
central extension of the loop algebra.
\end{proposition}
\begin{proof}
The central term is alternating by \(m=-n\) on its support.
In the cyclic sum of the three double brackets on \(x_m,y_n,z_s\),
the loop terms cancel by Jacobi in \(\mathfrak g\). A central term
occurs only for \(m+n+s=0\), and invariance of \(B\) identifies
the three cyclic values \(B([x,y],z)\), \(B([y,z],x)\),
and \(B([z,x],y)\). Their common coefficient is
\((m+n)+(n+s)+(s+m)=0\). This proves Jacobi on the extension.
The degree operator is a derivation: the bracket has loop
degree \(m+n\) and a central component can occur only in
degree zero. Adjoining \(d\) with this derivation proves
Jacobi for the stated algebra.
For \(h,h'\in\mathfrak h\) with \(B(h,h')\ne0\), the two
commuting loops \(h_1,h'_{-1}\) have bracket \(B(h,h')K\)
in the extension. Their lifts differ only by central
vectors, so a Lie section cannot exist.
\end{proof}

\begin{proposition}[Triangular decomposition and simple triples]
\label{res:LO10}
Set
\[
 \widehat{\mathfrak h}=\mathfrak h\oplus\mathbb C K\oplus\mathbb C d,
 \qquad
 \widehat{\mathfrak n}^{\pm}
 =\mathfrak n^\pm\otimes 1\ \oplus\
    \bigoplus_{\pm m>0}\mathfrak g\otimes t^m .
\]
These subspaces give a triangular decomposition.
Let \(\delta\) vanish on \(\mathfrak h,K\) and satisfy
\(\delta(d)=1\), and extend all finite roots by zero on \(K,d\).
The nonzero weights for the adjoint Cartan action are
\(\alpha+m\delta\) for finite roots \(\alpha\), and
\(m\delta\) for \(m\ne0\), with multiplicities respectively
\(1\) and \(\dim\mathfrak h\).
If \(\theta\) is the highest finite root, put
\[
 \alpha_0=\delta-\theta,\quad
 e_0=(f_\theta)_1,\quad
 f_0=(e_\theta)_{-1},\quad h_0=K-h_\theta .
\]
The finite simple-root triples and this extra triple satisfy
the rank-one triple relations and
\([e_i,f_j]=0\) for \(i\ne j\), including index \(0\).
The respective \(e_i\) generate \(\widehat{\mathfrak n}^+\),
and the \(f_i\) generate \(\widehat{\mathfrak n}^-\).
\end{proposition}
\begin{proof}
The degrees and finite Cartan eigenspaces give the decomposition
and the listed roots directly. The triangular summands are
Lie subalgebras: same-sign nonzero degrees add and cannot give
a central term, and degree zero uses the finite decomposition.
The highest root \(\theta\) is dominant by the root-string formula
and the absence of \(\theta+\alpha_i\).
Take any long root and make it dominant, say \(\beta\),
using the finite Weyl chamber result. Then
\(\theta-\beta\in Q_+\) by the highest-root bound, and
\((\theta-\beta,\theta+\beta)\ge0\) by dominance.
Thus \(\theta\) is long
\cite[Results D04, E03 and E09, PDF pp.~24--25, 30--31, 34--35]{mg-lie}.
Its coroot triple can be chosen with \([e_\theta,f_\theta]=h_\theta\).
Invariance gives
\(B(h_\theta,h)=\theta(h)B(e_\theta,f_\theta)\).
Since the root has squared length \(2\), its coroot \(h_\theta\)
is \(B\)-dual to \(\theta\), giving \(B(e_\theta,f_\theta)=1\).
The bracket from Proposition~\ref{res:LO09} now gives
\([e_0,f_0]=K-h_\theta\), and the two Cartan relations
follow from the \(\pm\theta\) weights. For a mixed pair
\([e_j,f_0]\) or \([e_0,f_j]\), there is no central term.
Its finite bracket vanishes because neither
\(\theta+\alpha_j\) nor \(-\theta-\alpha_j\) is a root, by
\cite[Result E09, PDF pp.~34--35]{mg-lie}.

The finite \(\mathfrak n^+\) is generated by its simple root
vectors by \cite[Results F01--F03, PDF pp.~36--38]{mg-lie}.
The adjoint module is irreducible, since its submodules are
ideals of the simple \(\mathfrak g\). Its lowest root vector
\(f_\theta\) is killed by \(\mathfrak n^-\), so PBW in the
reverse order shows that \(U(\mathfrak n^+)f_\theta=\mathfrak g\)
under the adjoint action. In degree one the subalgebra
generated by the \(e_i\) therefore contains
\(\mathfrak g\otimes t\). Successive brackets with this space
give all positive degrees, since \([\mathfrak g,\mathfrak g]
=\mathfrak g\). The same argument with the highest vector
gives the negative subalgebra.
\end{proof}

\section{Affine highest weights and level}

\begin{proposition}[Affine Verma modules]
\label{res:LO11}
For \(\Lambda\in\widehat{\mathfrak h}^{\,*}\), induce the
one-dimensional weight module from
\(\widehat{\mathfrak h}\oplus\widehat{\mathfrak n}^+\), with
positive root operators acting by zero. The resulting
\(\widehat M(\Lambda)\) is universal for a highest vector
of weight \(\Lambda\). It has finite-dimensional weight spaces,
one-dimensional top, and a unique simple highest-weight quotient
\(\widehat L(\Lambda)\). On both, \(K\) acts by
\(\Lambda(K)\), called the level.
\end{proposition}
\begin{proof}
Use PBW for the ordered basis of the affine Lie algebra,
with negative roots first \cite[Result A05, PDF pp.~8--9]{mg-lie}.
This identifies the induced module linearly with
\(U(\widehat{\mathfrak n}^-)\) acting freely on the top
vector, and its universal property follows by mapping that
vector to any given highest vector.

The roots \(\alpha_0,\alpha_1,\ldots,\alpha_r\) are linearly
independent, as evaluation on \(d\) singles out \(\alpha_0\).
Every positive root is a nonzero nonnegative combination
of them. For a positive-degree real root use
\(\alpha+m\delta=m\alpha_0+(m\theta+\alpha)\) and the
highest-root bound on \(-\alpha\).
For an imaginary root use \(m\delta=m\alpha_0+m\theta\).
At degree zero use the ordinary positive roots. Among roots
of bounded height there are finitely many: the coefficient
of \(\alpha_0\) bounds the loop degree, and each degree
has finite total root multiplicity. Thus a fixed height in
the difference from \(\Lambda\) allows just finitely many
negative-root PBW monomials. This proves the weight-space
assertions. Every proper submodule misses the top, and
the sum of all proper submodules still misses the top
by the same weight-space interpolation as in
Theorem~\ref{res:LO02}. Its quotient is simple.
The level assertion follows from centrality and the value
on the generating highest vector.
\end{proof}

\begin{theorem}[Integrability of simple affine highest weights]
\label{res:LO12}
Call a weight module integrable if every \(e_i,f_i\), for
\(0\le i\le r\), acts locally nilpotently. The simple
\(\widehat L(\Lambda)\) is integrable exactly when
\(\Lambda(h_i)\in\mathbb Z_{\ge0}\) for all these indices.
Its level is then a nonnegative integer. If
\(h_\theta=\sum_{i=1}^{r}a_i^\vee h_i\), the level is
\[
   \Lambda(K)=\Lambda(h_0)+
      \sum_{i=1}^{r}a_i^\vee\Lambda(h_i) .
\]
\end{theorem}
\begin{proof}
For a highest vector \(v\), the calculation in
Example~\ref{res:LO08} applies to each simple triple.
If \(f_i^N v=0\) with \(N\) minimal, applying \(e_i\)
shows \(N(\Lambda(h_i)-N+1) f_i^{N-1}v=0\).
Hence \(\Lambda(h_i)=N-1\), proving necessity.

Conversely, write \(m_i=\Lambda(h_i)\ge0\) and take the highest vector
\(v_\Lambda\) of \(\widehat M(\Lambda)\). The vector
\(f_i^{m_i+1}v_\Lambda\) is killed by \(e_i\)
by the same formula. It is killed by every other \(e_j\)
because \([e_j,f_i]=0\). These \(e_j\) generate
\(\widehat{\mathfrak n}^+\) by Proposition~\ref{res:LO10},
so the vector is highest at weight strictly below
\(\Lambda\). Its generated submodule stays below the original
top by PBW. The vector therefore vanishes in the simple
quotient.

It remains to verify local nilpotence on every vector.
The adjoint action of each \(e_i,f_i\) on the affine algebra
is locally nilpotent. On loop vectors its finite Lie action
is repeated bracketing with a fixed finite root vector,
which shifts finite weights in one root direction and
eventually gives zero. The loop degree changes uniformly
for index \(0\), and any central term is killed at the
next commutator. A bracket with \(d\) is a multiple
of the acting vector, and a further bracket is zero.
Adjoint local nilpotence passes to the enveloping algebra,
since the operator is a derivation and a product involves
only finitely many factors. For \(a=e_i\) or \(f_i\) the identity
\[
   a^N(uv)=
   \sum_{j=0}^N {N\choose j}(\operatorname{ad}a)^j(u)a^{N-j}v
\]
follows by induction. Its terms vanish for large \(N\),
since \(a\) is nilpotent on the top vector and
\(\operatorname{ad}a\) is locally nilpotent on \(u\).
Such vectors \(uv\) span \(\widehat L(\Lambda)\).
Finally the coroot of the highest finite root is a positive root
in the dual root system, so its simple-coroot coefficients \(a_i^\vee\)
are nonnegative integers by
\cite[Results E02 and E07, PDF pp.~29, 33--34]{mg-lie}.
Substituting \(K=h_0+h_\theta\) proves the level formula.
\end{proof}

\begin{proposition}[Nonzero level and evaluation modules]
\label{res:LO13}
A nonzero finite-dimensional module on which \(K\) acts as a scalar
has level zero. For \(a\in\mathbb C^\times\), substitution
\(x_m\mapsto a^m x\), \(K\mapsto0\), gives evaluation modules
of the affine algebra without \(d\). An evaluation module
obtained from a nontrivial \(\mathfrak g\)-action has no action
of \(d\) extending it to \(\widehat{\mathfrak g}\).
\end{proposition}
\begin{proof}
Take \(h,h'\in\mathfrak h\) with \(B(h,h')\ne0\).
On a module of dimension \(N>0\) the trace of
\([h_1,h'_{-1}]=B(h,h')K\) is zero. If the level is
\(\ell\), this gives \(N B(h,h')\ell=0\), so \(\ell=0\).
Evaluation respects the loop brackets by multiplication
of the powers of \(a\), and kills the central term.
If an operator for \(d\) existed, its bracket with
the action of \(x_0\) would vanish. Its bracket with
the action of \(x_1=a x_0\) would then also vanish,
but by \([d,x_1]=x_1\) that bracket equals \(a x_0\).
Since \(a\ne0\), all \(x_0\) would act by zero,
contrary to the nontrivial action.
\end{proof}

\begin{example}[The level-one basic highest weight]
\label{res:LO14}
Let \(\Lambda_0\) vanish on \(\mathfrak h,d\) and satisfy
\(\Lambda_0(K)=1\). Then \(\widehat L(\Lambda_0)\) is
integrable and infinite-dimensional. For a highest vector \(v\),
and any \(h\in\mathfrak h\) with \(B(h,h)\ne0\),
the vectors \(h_{-1}^{\,j}v\), \(j\ge0\), are nonzero
and lie in distinct \(d\)-eigenspaces.
\end{example}
\begin{proof}
The weight has value \(1\) on \(h_0=K-h_\theta\)
and \(0\) on each finite \(h_i\), so
Theorem~\ref{res:LO12} gives integrability.
The commutator \([h_1,h_{-1}]\) acts by the nonzero
scalar \(B(h,h)\) and \(h_1v=0\).
Commuting \(h_1\) across \(j\) copies of \(h_{-1}\) gives
\(h_1h_{-1}^{\,j}v=jB(h,h)h_{-1}^{\,j-1}v\).
Induction gives nonvanishing. The bracket
\([d,h_{-1}]=-h_{-1}\) gives eigenvalues \(-j\) on these
vectors, proving independence and infinite dimension.
\end{proof}

For the finite category and reciprocity, compare
\cite[Sections 15.1, 16.3 and 20.1--20.4]{etingof-lie} and
\cite[Chapters 1 and 3]{humphreys-o}. For the untwisted affine
model and integrable highest weights see
\cite[Sections 3, 4.1--4.7]{hernandez-affine} and
\cite[Chapters 6--7, 9 and 12]{kac}. Twisted affine algebras,
the general Kac--Moody character formula, and Kazhdan--Lusztig
polynomials are natural continuations of this chapter.
\chapter{Geometric Representations and Satake}

We prove geometric Satake for a complex split torus with rational
coefficients. We also compute the convolution square of the
minuscule intersection complex for \(\operatorname{GL}_2\).
These computations use the classical topology on complex points.
For the setting of general complex connected reductive groups, see
Baumann--Riche~\cite[Sections 6--9 and 13--14]{baumannriche}
and Achar~\cite[Chapter 9]{achar}.

Put \(O=\mathbb C[[t]]\) and \(K=\mathbb C((t))\).
For a complex algebraic group \(G\) the affine Grassmannian
\(\operatorname{Gr}_G\) has complex points \(G(K)/G(O)\).
Its bounded lattice models below are regarded with their reduced
structures. A sheaf on their complex points uses the ordinary
analytic topology. A sheaf on the union has finite support in the
sense that it is carried by one finite-dimensional closed stage.
We use the cohomological shift convention
\(H^a(F[b])=H^{a+b}(F)\).

\section{The affine Grassmannian of a torus}

\begin{lemma}[The reduced lattice stages of a split torus]
\label{res:GS01}
For \(T=(\mathbb G_{m,\mathbb C})^r\), let
\(\Lambda=X_*(T)=\mathbb Z^r\). The reduced affine Grassmannian
of \(T\) is the union of the discrete finite sets
\(\{-N,\ldots,N\}^r\), with the point labelled \(\lambda\)
represented by \(t^\lambda=(t^{\lambda_1},\ldots,t^{\lambda_r})\).
Every point is fixed by \(T(O)\).
\end{lemma}

\begin{proof}
Every member of \(K^\times\) has the unique form \(t^a u\)
with \(a\in\mathbb Z\) and \(u\in O^\times\).
This identifies the quotient and the stabilizers on complex points.
For one coordinate, the bounded lattice model consists of the
\(t\)-invariant subspaces of \(t^{-N}O/t^NO\), in all dimensions.
Invariance is a closed condition in each ordinary Grassmannian,
since the composite of the subbundle, the operator \(t\), and the
projection to the quotient bundle must vanish.
The finite space \(t^{-N}O/t^NO\) is a cyclic \(\mathbb C[t]\)-module
with relation \(t^{2N}=0\). A submodule lifts to an ideal of
\(\mathbb C[t]\) containing \((t^{2N})\), so is itself the image
of \((t^d)\) for one \(0\leq d\leq 2N\).
Thus there is precisely one complex point in each allowed dimension.
The product of these bounded models gives the torus stage.

For the reduced structure, a reduced affine complex algebra of
finite type is detected by its evaluations at complex points.
This is Lemma 1.4 of \emph{Topics in Field and Galois Theory}
\cite[\texttt{PV04}, PDF p.~5]{mg-fieldtopics}.
When there are finitely many such points the evaluation map to a
finite product of copies of \(\mathbb C\) is injective. For distinct
points choose a function separating them, subtract its
value at one point and normalize at the other. Products give
functions taking value \(1\) at one point and \(0\) at all the
others. On each affine chart the reduced scheme is therefore a
finite union of reduced points. This gives the claimed discrete
stages and their natural inclusions.
\end{proof}

The reduced qualification matters for functors evaluated on rings
with nilpotents. For example, over
\(\mathbb C[\epsilon]/(\epsilon^2)\), the unit
\(1+\epsilon t^{-1}\) reduces to \(1\) and does not belong
to the positive loop subgroup. Our sheaves use complex points
with the indicated topology throughout.

\begin{proposition}[The equivariant perverse category of a torus]
\label{res:GS02}
The category
\(\operatorname{Perv}_{T(O)}(\operatorname{Gr}_T,\mathbb Q)\)
of finite-support middle-perverse sheaves is equivalent to
the category of \(\Lambda\)-graded, finite-dimensional
\(\mathbb Q\)-vector spaces. A sheaf is sent to the family
\((V_\lambda)_{\lambda\in\Lambda}\) of its stalks in degree zero.
\end{proposition}

\begin{proof}
For a smooth stratum \(S\) of complex dimension \(d\), the
middle-perversity conditions are
\(H^a(i_S^*F)=0\) for \(a>-d\) and
\(H^a(i_S^!F)=0\) for \(a<-d\), where \(i_S\) includes the
stratum. For the disjoint points of Lemma~\ref{res:GS01}, both
restrictions are the ordinary restriction and \(d=0\).
This leaves a finite-dimensional vector space in degree zero
at each point. The finite-stage support condition is exactly
finite support of the family.

Equivariance on the perverse heart means an isomorphism between
pullbacks by the action and the projection, the identity at
the group identity, compatible with multiplication. At a point
fixed by \(T(O)\) both pullbacks are constant. At any finite jet
level the parameter space for the group is
\((\mathbb C^\times)^r\times \mathbb C^{r(N-1)}\) for truncation
modulo \(t^N\), hence connected. An isomorphism between constant
local systems on this parameter space is constant, and its
normalization at the identity forces the identity everywhere.
The unique isomorphisms are compatible under changing jet level.
Equivalently the product coordinates of \(T(O)\) are path-connected,
by paths in the nonzero constant terms and linear contractions of
the other coefficients. All linear maps between stalks respect
this equivariance. This computes the equivariant perverse heart,
without any assertion about equivariant derived Ext groups.
\end{proof}

\begin{proposition}[Convolution and symmetry for the torus]
\label{res:GS03}
For finite-support sheaves \(V,W\) identified as in
Proposition~\ref{res:GS02}, convolution has the formula
\[
 (V\star W)_\nu
   =\bigoplus_{\lambda+\mu=\nu} V_\lambda\otimes_{\mathbb Q} W_\mu .
\]
Its unit is \(\mathbb Q\) at \(0\), and its associativity is
the usual tensor rebracketing. Its geometric commutativity uses
the ordinary tensor flip.
\end{proposition}

\begin{proof}
The ordered convolution space is the twisted product
\(T(K)\times^{T(O)}\operatorname{Gr}_T\).
Use \((ga,x)\sim(g,ax)\) for \(a\in T(O)\).
Since \(T(O)\) acts trivially on the second factor, the labels
are precisely ordered pairs. The map to \(\operatorname{Gr}_T\)
multiplies the two successive modifications, hence takes
\((t^\lambda,t^\mu)\) to \(t^{\lambda+\mu}\).
The descended exterior tensor product at the pair is
\(V_\lambda\otimes W_\mu\) in degree zero by the explicit
equivariance in Proposition~\ref{res:GS02}.
Pushforward on the finite support is direct sum over a finite
discrete fibre, with no higher cohomology. The same argument
for triples proves the associativity and unit formulas.

The two-point exchange can be computed over the affine line
with coordinate \(z\). For marked points \(a,b\), in each line
coordinate the ordered modifications are multiplication by
\((z-a)^{\lambda_i}\), followed by
\((z-b)^{\mu_i}\), in the punctured trivialization.
If \(a\ne b\) the factors give the independent orders at the two
points. If \(a=b\) they give the summed order in the resulting
modification. The ordered intermediate modification still records
the first exponent. Thus for fixed ordered labels these
modifications form a section over the whole two-point parameter
plane, with the constant coefficient space
\(V_\lambda\otimes W_\mu\). There is no monodromy on that family.
Exchanging the marked points and the labels interchanges the two
coefficient factors by the ordinary flip, also at the collision
point. Summing over labels gives the commutativity constraint
from fusion. All fibre degrees here are zero.
The symmetry identities are exactly the identities for flipping
and rebracketing ordinary tensor products.
\end{proof}

\begin{theorem}[Geometric Satake for a split torus]
\label{res:GS04}
Let \(\check T=\operatorname{Spec}\mathbb Q[\Lambda]\), the
split algebraic \(\mathbb Q\)-torus with character lattice
\(\Lambda=X_*(T)\). There is a symmetric monoidal equivalence
\[
 \bigl(\operatorname{Perv}_{T(O)}(\operatorname{Gr}_T,\mathbb Q),
          \star\bigr)
  \simeq
  \bigl(\operatorname{Rep}_{\mathbb Q}^{\mathrm{fd}}(\check T),
          \otimes_{\mathbb Q}\bigr),
\]
where the right side is the category of finite-dimensional rational
algebraic representations. The stalk at \(\lambda\) is sent
to the \(\lambda\)-weight space. Global cohomology on the left
is the forgetful fibre functor on the right and is concentrated
in degree zero.
\end{theorem}

\begin{proof}
Write \(e^\lambda\) for the basis element of the group algebra.
Its comultiplication and counit are
\(\Delta(e^\lambda)=e^\lambda\otimes e^\lambda\) and
\(\varepsilon(e^\lambda)=1\). An algebraic representation is a
finite-dimensional comodule for this group algebra. The coaction
calculation in Proposition~\ref{res:IN03}, applied to the algebraic
coaction over \(\mathbb Q\), identifies such a comodule with its
direct sum of weight spaces.
Conversely, a finite grading gives the coaction
\(v\mapsto v\otimes e^\lambda\) on each homogeneous vector of
weight \(\lambda\), and comodule maps are precisely the maps of
graded vector spaces.
Under tensor product the coaction multiplies the two basis
elements, \(e^\lambda e^\mu=e^{\lambda+\mu}\).
The category equivalence of Proposition~\ref{res:GS02} now
identifies the two tensor rules and symmetries by
Proposition~\ref{res:GS03}. Global sections on finitely many
points are the direct sum of stalks and are exact.
\end{proof}

\begin{corollary}[Reconstruction of the dual torus]
\label{res:GS05}
For any commutative \(\mathbb Q\)-algebra \(A\), tensor
automorphisms of the fibre functor of Theorem~\ref{res:GS04}
after extension to \(A\) form the group
\(\operatorname{Hom}(\Lambda,A^\times)=\check T(A)\),
naturally in \(A\).
\end{corollary}

\begin{proof}
The one-dimensional object supported at \(\lambda\) has
base-extended fibre \(A\). Its tensor automorphism is given
by some \(u_\lambda\in A^\times\). For an arbitrary object
every rational vector of weight \(\lambda\) defines a map from
the corresponding one-dimensional object. Naturality forces
multiplication by \(u_\lambda\) on that whole extended weight
space. The unit and convolution rule give
\(u_0=1\) and \(u_{\lambda+\mu}=u_\lambda u_\mu\).
Conversely these conditions define a natural tensor automorphism.
An algebra map \(\mathbb Q[\Lambda]\to A\) is exactly the
assignment \(e^\lambda\mapsto u_\lambda\) with these conditions.
\end{proof}

\section{Local sheaf calculations}

We use the derived category, ordinary truncation, and triangle
conventions of \emph{Advanced Homological Algebra}
\cite[Results D03, D05 and D08, PDF pp.~24--27]{mg-ha}.
The following lemma supplies the comparison of singular and sheaf
cohomology needed for the explicit lattice charts.

\begin{lemma}[Constant sheaves and small singular cochains]
\label{res:GS06}
On a second-countable locally compact Hausdorff space with
a basis of contractible open sets at every point,
cohomology of the constant rational sheaf agrees naturally,
under restriction to open sets, with ordinary singular
cohomology. For a closed subset \(Z\) of an open set \(U\),
cohomology with support in \(Z\) is the singular cohomology
of the pair \((U,U\setminus Z)\).
\end{lemma}

\begin{proof}
Sheafify the presheaves of rational singular cochains
\(C^q(U;\mathbb Q)\) to obtain sheaves \(\mathcal S^q\).
Filtered direct limits of vector spaces are exact.
At each stalk, a basis of contractible opens therefore
gives an exact augmented complex
\(\mathbb Q_X\longrightarrow\mathcal S^\bullet\).

We verify that the cochain map to sections of each
\(\mathcal S^q\) over any open \(U\) is onto.
A section is represented locally by cochains \(c_i\) over
opens \(V_i\). Choose a locally finite covering by smaller
opens \(W_i\) with \(\overline{W_i}\subset V_i\), taking closures
in \(U\). Such a refinement is obtained from a compact
exhaustion of \(U\) whose terms lie in the interiors of
the next terms. Cover each compact annulus by finitely many
relatively compact opens inside representing opens,
with each chosen open avoiding all but a bounded number
of adjoining annuli.
For a simplex with image in one or more \(W_i\), give it
the value of one corresponding \(c_i\). Give all other
simplices value zero. At any \(x\in U\), choose \(W_j\)
containing \(x\). Local finiteness lets us shrink a
neighborhood of \(x\) to meet only those \(W_i\) whose
closures contain \(x\). There are finitely many. For these
indices the germs of \(c_i\) and \(c_j\) agree at \(x\),
so they agree after restriction to a still smaller
neighborhood. On simplices in that neighborhood all
choices give the same value. The constructed cochain thus
represents the section. Extending a cochain by zero on
simplices not wholly contained in its original open set
proves that \(\mathcal S^q\) is flabby, meaning all open
restriction maps are onto.

We recall why flabby resolutions compute the derived
sections used here. A sheaf of rational vector spaces embeds
in the product over points of the skyscraper sheaves of
its stalks. Each skyscraper of a vector space is injective,
since maps into it are linear maps from a stalk and the
stalk functor is exact. A product of these sheaves is
injective and flabby. If \(F\) is flabby in a short exact
sequence \(0\to F\to G\to H\to0\), every section of \(H\)
on an open set lifts to \(G\) there. Indeed local lifts
exist, and their differences on overlaps are sections of
\(F\). Extend a difference using flabbiness to adjust each
new lift so it agrees with an existing one. A maximal
compatible lift exists by the sheaf axiom along chains
of partial lifts, and this adjustment extends it at any
missing point. It follows that an injective resolution of
a flabby sheaf constructed by successive products of
skyscrapers has flabby successive quotients and stays exact
on sections. This proves the required acyclicity.

The kernel of \(C^\bullet(U;\mathbb Q)\to
\Gamma(U,\mathcal S^\bullet)\) consists of the cochains that
vanish on all simplices in some open covering, each
simplex required to lie in one member of the covering.
For a fixed covering this is the dual of the quotient by
its small chains. That quotient is acyclic by the
small-chain comparison of \emph{Topology and Riemannian
Geometry}~\cite[Theorem 3.8, PDF p.~19]{mg-topriem}.
Dualizing vector spaces and taking the filtered union over
covering refinements preserve acyclicity. This gives the
absolute comparison on each open set. The flabby resolution
has an onto restriction to \(U\setminus Z\) in every degree,
and the kernel computes sections with support in \(Z\)
by comparison with injective resolutions and the same
restriction sequence. The singular cochain restriction is
also onto, by extension by zero. Comparing the two kernels
via their long exact sequences gives the relative comparison.
\end{proof}

\paragraph{Stalks and costalks.}
For an open immersion \(j:U\hookrightarrow X\), the stalks
of \(Rj_*\mathbb Q_U\) are the direct limits of cohomology of
intersections of neighborhoods with \(U\). More generally,
for \(m:Y\to X\), the stalks of \(Rm_*\mathbb Q_Y\) are the
direct limits over neighborhoods \(B\) of the cohomology of
\(m^{-1}(B)\). In each assertion one applies the direct image
to an injective resolution, then takes the stalk, which is
an exact filtered direct limit. The costalk of \(\mathbb Q_X\)
at a closed point is computed by the corresponding direct
limit with support at the point. Indeed the right adjoint
to extension from a closed point selects sections supported
there, and applying it to an injective resolution gives
this formula. We use Lemma~\ref{res:GS06} in these expressions.

\section{Two lattice steps for \texorpdfstring{\(\operatorname{GL}_2\)}{GL(2)}}

Let \(V=\mathbb C^2\) and \(L_0=O^2\). A lattice in \(K^2\)
is a free rank-two \(O\)-submodule spanning \(K^2\).
A matrix in \(\operatorname{GL}_2(K)\) carries \(L_0\) to
the lattice generated by its columns, and the stabilizer
of \(L_0\) is \(\operatorname{GL}_2(O)\). Thus the quotient
has these lattices as its points.

\begin{proposition}[The minuscule orbit and the two-step Schubert variety]
\label{res:GS07}
The lattices \(L\subset L_0\) with
\(\dim_{\mathbb C}(L_0/L)=1\) form the orbit
\(X_{(1,0)}\simeq\mathbb P(V)\).
Those with \(\dim_{\mathbb C}(L_0/L)=2\) form the closed
Schubert variety \(X_{(2,0)}\), with one orbit of type
\((2,0)\) and the closed orbit \(X_{(1,1)}=\{tL_0\}\).
\end{proposition}

\begin{proof}
On each of these finite-dimensional quotients the
operator \(t\) is nilpotent because a nonzero eigenvalue
\(a\) would contradict invertibility of \(t-a\) in \(O\).
For colength one this means \(tL_0\subset L\).
The quotient \(L/tL_0\) is a line in \(V\), and every line
gives its inverse-image lattice. This identifies the
bounded algebraic model, and constant matrices act
transitively.

For colength two we have \(t^2L_0\subset L\).
Modulo \(t^2L_0\) the lattices are the two-dimensional
subspaces of \(W=L_0/t^2L_0=V\oplus tV\) stable under
\(N(x,y)=(0,x)\). They form a closed subset of an ordinary
Grassmannian by the same subbundle criterion as in
Lemma~\ref{res:GS01}. If \(t\) induces zero on \(L_0/L\),
then \(L=tL_0\). Otherwise a vector and its nonzero \(t\)-image
give a complex basis of the quotient, making \(L_0/L\) cyclic
over \(O/(t^2)\). The resulting
surjection \(L_0\to O/(t^2)\) has a row of coefficients
one of which is a unit modulo \(t^2\).
After scaling the corresponding basis vector and subtracting
its multiples from the other, the kernel is
\(t^2O\oplus O\). Hence this locus is one orbit and is
open because its complement is the single lattice \(tL_0\).
To see that the orbit closure contains that lattice,
fix a nonzero endomorphism \(B\) of \(V\) with \(B^2=0\).
The graphs \(\{(sBy,y):y\in V\}\) are \(N\)-stable, give
quotients of the nonzero-\(t\) type for \(s\ne0\), and tend
to \(tV\) at \(s=0\).
This proves the asserted closure and both orbit descriptions.
\end{proof}

\begin{proposition}[The convolution surface and its fibres]
\label{res:GS08}
The convolution of the two colength-one lattice stages has
ordered space
\[
 \widetilde X=\{L_0\supset L_1\supset L_2:
       \dim_{\mathbb C}(L_0/L_1)=\dim_{\mathbb C}(L_1/L_2)=1\}.
\]
It is the projective bundle of lines in
\(\mathcal O(-1)\oplus\mathcal O(1)\) over \(\mathbb P^1\).
The map \(m:\widetilde X\to X_{(2,0)}\) retaining \(L_2\)
is an isomorphism over the orbit of type \((2,0)\).
Its fibre over \(tL_0\) is a section isomorphic to
\(\mathbb P^1\) with normal line bundle \(\mathcal O(-2)\).
\end{proposition}

\begin{proof}
A point in the twisted convolution product specifies its
first lattice \(L_1\) and a colength-one modification
relative to \(L_1\), namely \(L_2\).
Changing a basis of \(L_1\) gives exactly the
\(\operatorname{GL}_2(O)\)-identification in the twisted
product. A choice of the second lattice amounts to a line
in \(L_1/tL_1\).
Let \(\mathcal S\subset V\otimes\mathcal O\) be the
tautological line on the first \(\mathbb P(V)\).
In the truncated power-series coordinates one has
\[
 L_1/tL_1=\mathcal S\oplus t(V/\mathcal S).
\]
Indeed \(L_1\) has constant term in \(\mathcal S\) and
arbitrary positive coefficients, while \(tL_1\) has
first coefficient in \(\mathcal S\) and unrestricted
coefficients of order two or higher. The quotient line
\(V/\mathcal S\) is \(\mathcal S^\vee\) after fixing a
volume form on \(V\), since its pairing with \(\mathcal S\)
by wedge is perfect. This gives the projective bundle,
which is smooth and projective. Equivalently the two
lattice conditions are closed incidence conditions in
finite flag varieties, so \(m\) is proper.

If \(L_0/L_2\simeq O/(t^2)\), it has a unique length-one
submodule, the image of multiplication by \(t\).
Thus \(L_1=L_2+tL_0\). The sum is an algebraic function
on this rank-constant open locus, since the image of a
constant-rank matrix of bundle maps is locally a
subbundle. This gives the inverse over the open orbit.
Over \(L_2=tL_0\) any first line is permitted. The line
in \(L_1/tL_1\) at this section is \(t(V/\mathcal S)\).
Locally neighboring lines are graphs of maps to \(\mathcal S\),
so the normal bundle is
\(\operatorname{Hom}(t(V/\mathcal S),\mathcal S)=\mathcal O(-2)\).
In particular the unique
positive-dimensional fibre has twice its dimension equal
to the codimension of its image stratum.
\end{proof}

\begin{lemma}[The transverse cone at the closed orbit]
\label{res:GS09}
A neighborhood of \(x=tL_0\) in \(X_{(2,0)}\) is the
vertex neighborhood in the quadratic cone
\(\{(u,v,w)\in\mathbb C^3:u^2+vw=0\}\).
In the classical topology it is locally the cone on
\(\mathbb{RP}^3\). The inverse image of the affine cone
chart in \(\widetilde X\) is the total space of
\(\mathcal O(-2)\) over \(\mathbb P^1\), with the zero
section mapping to \(x\).
\end{lemma}

\begin{proof}
Use the chart of the Grassmannian of two-planes in \(W\)
where projection to the second copy of \(V\) is invertible.
There a plane has the form \(\{(Ay,y):y\in V\}\) for
\(A\in\operatorname{End}_{\mathbb C}(V)\). It is stable
under \(N\) precisely when \(A^2=0\).
Set-theoretically this is equivalent to zero trace and
zero determinant. One direction follows from nilpotence,
and the converse follows from the characteristic identity
for a \(2\) by \(2\) matrix.
Writing \(A=\left(\begin{smallmatrix}u&v\\w&-u\end{smallmatrix}\right)\)
gives the displayed equation in the reduced chart.

The map \((a,b)\mapsto(ab,-a^2,b^2)\) from \(\mathbb C^2\)
is onto this cone. If \(v\ne0\), take \(a^2=-v\) and
\(b=u/a\). If \(v=0\), then \(u=0\) and one takes \(a=0\).
Each fibre is precisely an orbit of the simultaneous sign
change, including the singleton fibre at \(0\).
The equalities \(|a|^2=|v|\) and \(|b|^2=|w|\) show the
map is proper. Hence the induced continuous bijection
from the quotient is a homeomorphism. Quotienting radial
coordinates identifies the link with \(S^3/\{\pm1\}
=\mathbb{RP}^3\).

On this chart a first lattice is specified by a line
\(\ell\subset V\). The inclusion \(L_2\subset L_1\) is
\(\operatorname{im}A\subset\ell\), and the inclusion
\(tL_1\subset L_2\) is \(\ell\subset\ker A\).
These are precisely the graphs described by a linear
map \(V/\ell\to\ell\). They form the total space of the
normal line bundle from Proposition~\ref{res:GS08}.
For a small ball about \(A=0\) in the graph chart its
inverse image retracts onto the zero section by
\((\ell,A)\mapsto(\ell,sA)\), \(0\leq s\leq1\).
\end{proof}

\begin{theorem}[The intersection complex of the two-step orbit closure]
\label{res:GS10}
Over \(\mathbb Q\), the intersection complex of \(X=X_{(2,0)}\),
normalized to be perverse, is \(\mathbb Q_X[2]\).
At \(x=tL_0\) it has one-dimensional stalk in degree \(-2\)
and one-dimensional costalk in degree \(2\), with the other
stalk and costalk degrees zero.
\end{theorem}

\begin{proof}
By Proposition~\ref{res:GS08}, the complement \(U=X\setminus\{x\}\)
is a smooth complex surface. By Lemma~\ref{res:GS09}, sufficiently
small contractible cone neighborhoods \(B\) of \(x\) have
punctured neighborhoods retracting to \(\mathbb{RP}^3\).
The rational cohomology of that link is \(\mathbb Q\) in
degrees \(0,3\) only, by the projective-space calculation and
universal coefficients in \emph{Topology and Riemannian Geometry}
\cite[Proposition 4.3 and Theorem 4.13,
PDF pp.~24 and 29]{mg-topriem}. The degree-zero restriction
from \(B\) to its punctured neighborhood is an isomorphism.
The relative cohomology of \((B,B\setminus\{x\})\) is therefore
\(\mathbb Q\) in degree \(4\) only by the pair sequence
\cite[Proposition 4.8, PDF p.~26]{mg-topriem}.
These identifications are compatible on a nested radial basis
by radial deformation. Lemma~\ref{res:GS06} computes the
stalk and costalk of the shifted constant sheaf as asserted.
On the open surface the only degree is \(-2\), and the
closed point satisfies the degree-zero bounds. These are
the two checks for perversity.

For this single closed stratum, the middle-perversity
intersection complex is constructed as
\(\tau_{\leq-1}Rj_*\mathbb Q_U[2]\) with \(j:U\hookrightarrow X\)
and ordinary truncation \(\tau\).
Before the shift by \(2\), the middle-perversity cutoff in real
codimension four is degree \(1\). These conventions are given by
de Cataldo--Migliorini~\cite[Sections 2.1 and 2.7]{decataldomigliorini}.
On \(U\) the natural map from \(\mathbb Q_X[2]\) to this
truncation is an isomorphism. At \(x\), the untruncated
target has degrees \(-2,1\) by the punctured-neighborhood
calculation, and the map is an isomorphism in degree \(-2\)
by connectedness. The truncated map is thus also a
stalkwise quasi-isomorphism at \(x\).
Adjunction computes morphisms from and to a sheaf supported
at \(x\) using the costalk and stalk, respectively. The strict
bounds make their degree-zero Hom groups vanish. This verifies
the middle-extension property on the perverse heart.
\end{proof}

\begin{theorem}[The minuscule convolution square for \(\operatorname{GL}_2\)]
\label{res:GS11}
Write \(\operatorname{IC}_{(a,b)}\) for the
rational middle-perversity intersection complex of the
indicated orbit closure, extended by zero to the affine
Grassmannian. For \(G=\operatorname{GL}_{2,\mathbb C}\),
\[
 \operatorname{IC}_{(1,0)}\star\operatorname{IC}_{(1,0)}
       \simeq \operatorname{IC}_{(2,0)}
                   \oplus\operatorname{IC}_{(1,1)}
\]
as \(G(O)\)-equivariant perverse sheaves.
\end{theorem}

\begin{proof}
The minuscule factor is \(\mathbb Q_{\mathbb P^1}[1]\).
In local bases of the first lattice both factors are constant
with trivial actions on their coefficients. Thus the twisted
exterior square is \(\mathbb Q_{\widetilde X}[2]\) and the convolution is
\(Rm_*\mathbb Q_{\widetilde X}[2]\). By
Proposition~\ref{res:GS08}, over \(U\) the unit of pullback
and pushforward
\[
 \mathbb Q_X[2]\longrightarrow Rm_*\mathbb Q_{\widetilde X}[2]
\]
is an isomorphism. At \(x\), the inverse images of
small radial neighborhoods retract to the zero section
\(\mathbb P^1\) by Lemma~\ref{res:GS09}. The derived
pushforward has stalk cohomology \(\mathbb Q\) in degrees
\(-2,0\) and zero otherwise. Here we used
Lemma~\ref{res:GS06} and the complex projective-space
cohomology from \emph{Topology and Riemannian Geometry}
\cite[Proposition 4.3 and Theorem 4.13,
PDF pp.~24 and 29]{mg-topriem}. The unit identifies degree
\(-2\) on the two sides. Its cone has just one cohomology
sheaf, a degree-zero copy of \(i_*\mathbb Q\) for the
inclusion \(i:\{x\}\hookrightarrow X\).

The connecting morphism of this triangle belongs to
\[
 \operatorname{Hom}(i_*\mathbb Q,\mathbb Q_X[3])
   = H^3(i^!\mathbb Q_X)=0
\]
by the costalk computation of Theorem~\ref{res:GS10}.
Therefore the identity of the cone has a lift to a section
of the quotient map, by applying \(\operatorname{Hom}(i_*\mathbb Q,-)\)
to the triangle. The original map and that section together
give an isomorphism, as they are an isomorphism on
cohomology sheaves. The section is unique since
\(\operatorname{Hom}(i_*\mathbb Q,\mathbb Q_X[2])=0\) as well.
The two summands are perverse by Theorem~\ref{res:GS10}
and the bound at a point, respectively.

For equivariance the action on the finite lattice diagrams
factors through the smooth group \(H=\operatorname{GL}_2(O/t^2O)\).
The pullback of the displayed triangle under the projection
from \(H\times X\) is the same triangle for
\(\operatorname{id}_H\times m\). For the pushforward term
the natural comparison is a stalkwise isomorphism because a
small contractible neighborhood in \(H\) times \(m^{-1}(B)\)
has the cohomology of \(m^{-1}(B)\).
The action pullback has the same description after
the isomorphism \((h,y)\mapsto(h,hy)\).
These identifications give the normalized equivariance
and preserve the unit. The cone has trivial coefficient
action at \(x\), since each \(h\) preserves the complex
orientation of the projective-line fibre and thus its
top cohomology. Under pullback to \(H\times X\), the
same uniqueness argument still applies. Indeed the
costalk sheaf along \(H\times\{x\}\) has cohomology
sheaves only in degree \(4\) before shifting. Locally a
contractible neighborhood of \(h\) times
the radial pair has the same relative cohomology
as the radial pair itself, by Lemma~\ref{res:GS06}.
Its hypercohomology therefore vanishes below degree \(4\),
which gives both Hom vanishings used for uniqueness.
The two pullbacks of the splitting agree with the
equivariance isomorphism. This proves the asserted
equivariant decomposition.
\end{proof}

\begin{corollary}[Central translates of the minuscule rule]
\label{res:GS12}
For integers \(a,b\) and a dominant pair \(\lambda_1\geq\lambda_2\),
convolution with the intersection complex of the point \(t^aL_0\) sends
\(\operatorname{IC}_{(\lambda_1,\lambda_2)}\) to
\(\operatorname{IC}_{(\lambda_1+a,\lambda_2+a)}\).
In particular,
\[
 \operatorname{IC}_{(a+1,a)}\star\operatorname{IC}_{(b+1,b)}
 \simeq \operatorname{IC}_{(a+b+2,a+b)}
             \oplus\operatorname{IC}_{(a+b+1,a+b+1)} .
\]
\end{corollary}

\begin{proof}
The scalar \(t^a\) is central in the loop group and commutes
with the \(G(O)\)-action. The convolution map with that
point in either position is the isomorphism of lattice
spaces \(L\mapsto t^aL\). The twisted tensor factor at the
point is \(\mathbb Q\) in degree zero, so the pushforward
is just transport by that isomorphism, preserving the
middle-perversity normalization and the orbit labels
up to addition of \((a,a)\).
For the last formula, translating the first lattice by
\(t^a\) and the final lattice by \(t^{a+b}\) identifies
the two translated steps with the ordered minuscule
steps whose terminal lattice is translated by \(t^{a+b}\).
Apply Theorem~\ref{res:GS11} and the first formula.
\end{proof}
\backmatter
\begingroup
\footnotesize
\begin{thebibliography}{99}
\phantomsection
\addcontentsline{toc}{chapter}{Bibliography}
\renewcommand{\runningchapter}{Bibliography}
\setlength{\itemsep}{2pt}

\bibitem{achar}
P.~N. Achar.
\emph{Perverse Sheaves and Applications to Representation Theory}.
Mathematical Surveys and Monographs 258, American Mathematical Society, 2021.

\bibitem{assem}
I. Assem, D. Simson, and A. Skowro\'nski.
\emph{Elements of the Representation Theory of Associative Algebras}, Vol.~1.
Cambridge University Press, 2006.

\bibitem{baumannriche}
P. Baumann and S. Riche.
\emph{Notes on the Geometric Satake Equivalence}.
arXiv:1703.07288v3, 2018.

\bibitem{decataldomigliorini}
M.~A.~A. de Cataldo and L. Migliorini.
\emph{The Decomposition Theorem, Perverse Sheaves and the Topology of Algebraic Maps}.
arXiv:0712.0349v2.

\bibitem{derksen}
H. Derksen and G. Kemper.
\emph{Computational Invariant Theory}.
Second enlarged edition, Springer, 2015.

\bibitem{etingof-rep}
P. Etingof, O. Golberg, S. Hensel, T. Liu, A. Schwendner,
D. Vaintrob, and E. Yudovina.
\emph{Introduction to Representation Theory}.
MIT 18.712 course edition, 2011.

\bibitem{etingof-lie}
P. Etingof.
\emph{Representations of Lie Groups}.
MIT 18.757 course notes, Fall 2023.

\bibitem{hernandez-affine}
D. Hernandez.
\emph{An Introduction to Affine Kac--Moody Algebras}.
CTQM Master Class Series 2, 2006.

\bibitem{humphreys-o}
J.~E. Humphreys.
\emph{Representations of Semisimple Lie Algebras in the BGG Category O}.
Graduate Studies in Mathematics 94, American Mathematical Society, 2008.

\bibitem{kac}
V.~G. Kac.
\emph{Infinite-Dimensional Lie Algebras}.
Third edition, Cambridge University Press, 1990, corrected reprint, 1995.

\bibitem{kraft-procesi}
H. Kraft and C. Procesi.
\emph{Classical Invariant Theory: A Primer}.
Lecture notes, 1996.

\bibitem{milneag}
J.~S. Milne.
\emph{Algebraic Groups: The Theory of Group Schemes of Finite Type over a Field}.
Cambridge University Press, corrected reprint, 2022.

\bibitem{serre-finrep}
J.-P. Serre.
\emph{Linear Representations of Finite Groups}.
Translated by L.~L. Scott, Graduate Texts in Mathematics 42, Springer, 1977.

\bibitem{vdk}
W. van der Kallen.
\emph{Reductivity Properties over an Affine Base}.
arXiv preprint, version 2.

\bibitem{webb}
P. Webb.
\emph{A Course in Finite Group Representation Theory}.
Author's prepublication manuscript, February 2016.

\bibitem{mg-algebra-core}
MathGenome.
\emph{Algebra Core}. P1-03, revision in MathGenome v0.25.0.

\bibitem{mg-alggroup}
MathGenome.
\emph{Linear Algebraic Groups}. ALG-09, 55-page revision, 2026.

\bibitem{mg-d07}
MathGenome.
\emph{Abstract Algebra}. D07, revision in MathGenome v0.25.0.

\bibitem{mg-fieldtopics}
MathGenome.
\emph{Topics in Field and Galois Theory}. ALG-04, 31-page revision, 2026.

\bibitem{mg-ha}
MathGenome.
\emph{Advanced Homological Algebra}. ALG-10, 62-page revision in
MathGenome v0.25.0.

\bibitem{mg-lie}
MathGenome.
\emph{Lie Algebras and Finite-Dimensional Representations}.
ALG-08, 74-page revision, 2026.

\bibitem{mg-rm03}
MathGenome.
\emph{Rings, Modules, and Associative Algebras}. ALG-03, 19-page core notes, 2026.

\bibitem{mg-rm03d}
MathGenome.
\emph{Topics in Rings, Modules, and Associative Algebras}.
ALG-03, 38-page deepening notes, 2026.

\bibitem{mg-topriem}
MathGenome.
\emph{Topology and Riemannian Geometry}. P2-02, revision in MathGenome v0.25.0.

\end{thebibliography}
\endgroup
\end{document}
